\documentclass{book}

\usepackage[utf8]{inputenc}  % para que funcionen las tildes
\usepackage{amsmath}
\usepackage{amssymb}
\usepackage{amsthm}
\usepackage{cancel}
\usepackage{letltxmacro}
\LetLtxMacro\amsproof\proof
\LetLtxMacro\amsendproof\endproof
\usepackage{txfonts}
\usepackage{lmodern}
\usepackage[dvipsnames]{xcolor}
\usepackage{thmtools, thm-restate}
\usepackage{datetime} % para la hora de compilación
\usepackage{graphicx}
\graphicspath{{images/}}
\usepackage[spanish,es-noquoting]{babel} % es-noquoting es para que funcione tikz
\usepackage{mathabx} % para \divides
\usepackage{centernot} % para \centernot\inculdes
\usepackage{wrapfig}
\usepackage{multicol}
\usepackage{subcaption}
\usepackage{etoolbox}
\usepackage{xfrac}
\usepackage{glossaries}
\usepackage{standalone}
\usepackage{pgfmath}
\usepackage{pgfplots}
\usepackage{pdfpages}
\usepackage{tikz}
\usepackage{mathtools}
\usepackage{hyperref}
\usepackage{cancel}
\usepackage[shortlabels]{enumitem}
\usetikzlibrary{arrows.meta}
\usetikzlibrary{shapes}
\usetikzlibrary{decorations.text}
\usetikzlibrary{positioning}
\usetikzlibrary{external}
\tikzexternalize[prefix=./tikzbuild/]
\tikzexternalize % activate!

\AtBeginDocument{%
  \LetLtxMacro\proof\amsproof
  \LetLtxMacro\endproof\amsendproof
}

\definecolor{coolblack}{rgb}{0.0, 0.18, 0.39}
%\usepackage[a6paper,margin=5mm]{geometry}
\usepackage[a4paper, top=2.5cm,bottom=2.5cm,left=2.5cm,right=2.5cm]{geometry}

%\setlength{\parindent}{0pt}
\usepackage{parskip}

\usepackage{hyperref}

% ESTILOS DE DEFINICIONES Y TEOREMAS
\declaretheoremstyle[
	bodyfont=\normalfont,
	shaded={
		margin=8pt,
		bgcolor=White,
		rulecolor=Black,
		rulewidth=1pt
}]{mythm}

\declaretheoremstyle[
	bodyfont=\normalfont,
	shaded={
		margin=1em,
		bgcolor={rgb}{0.9,0.9,0.9}
}]{mydfn}

\declaretheoremstyle[
bodyfont=\normalfont,
spacebelow=1em,
spaceabove=1em,
]{myej}

\declaretheoremstyle[
bodyfont=\normalfont,
spacebelow=1em,
spaceabove=1em,
shaded={
    margin=8pt,
    bgcolor=White,
    rulecolor=coolblack,
    rulewidth=1pt
}]{thej}

\declaretheoremstyle[
bodyfont=\normalfont,
spacebelow=1em,
spaceabove=1em,
thmbox=M
]{myeg}

% DEFINICIONES DE ENTORNOS DE TEOREMAS
\declaretheorem[
	name=Teorema,
	refname={teorema,teoremas},
	Refname={Teorema,Teoremas},
	style=mythm
]{thm}

\declaretheorem[
	name=Corolario,
	refname={corolario,corolarios},
	Refname={Corolario,Corolarios},
	style=myej
]{cor}

\declaretheorem[
	name=Proposici\'{o}n,
	refname={proposici\'{o}n,proposiciones},
	Refname={Proposici\'{o}n, Proposiciones},
	sharenumber=thm,
	style=myej
]{pro}

\declaretheorem[
name=Lema,
refname={lema,lemas},
Refname={Lema, Lemas},
sharenumber=thm,
style=myej
]{lm}

%\newtheorem*{cor}{Corolario}
\newtheorem*{lem}{Lema}

\declaretheorem[
	name=Definici\'{o}n,
	refname={definici\'{o}n,definiciones},
	Refname={Definici\'{o}n,Definiciones},
	style=mydfn
]{dfn}

\declaretheorem[
name=Ejercicio,
refname={ejercicio,ejercicios},
Refname={Ejercicio,Ejercicios},
style=myej,
numbered=no
]{ex}

\declaretheorem[
name=Ejercicio propuesto,
refname={ejercicio,ejercicios},
Refname={Ejercicio,Ejercicios},
style=thej,
]{th_ex}

\declaretheorem[
name=Ejemplo,
refname={ejemplo,ejemplos},
Refname={Ejemplo,Ejemplos},
style=myeg
]{eg}

\declaretheorem[
name=Observaci\'{o}n,
refname={observaci\'{o}n,observaciones},
Refname={Observaci\'{o}n,Observaciones},
style=myej,
numbered=no
]{obs}

\makeglossaries
\newglossaryentry{di}{name={DI},description={Dominio de integridad (ver \autoref{dfn:dominiointegridad})}}
\newglossaryentry{dip}{name={DIP},description={Dominio de ideales principales (ver \autoref{dfn:dominioidealesprincipales})}}
\newglossaryentry{dfu}{name={DFU},description={Dominio de factorización única (ver \autoref{dfn:dfu})}}
\newglossaryentry{de}{name={DE},description={Dominio euclídeo (ver \autoref{dfn:de})}}


%\setcounter{chapter}{-1}
%COMANDOS ÚTILES PARA DERIVADAS
\AtBeginDocument{\renewcommand{\d}{\mathrm{d}}}
\newcommand{\Dd}[1]{\frac{\mathrm{d}}{\mathrm{d}#1}}
\newcommand{\dd}[2]{\frac{\mathrm{d}#1}{\mathrm{d}#2}}
\newcommand{\mbf}[1]{\mathbf{#1}}
% COMANDOS ÚTILES PARA LA TEORÍA DE GRUPOS
%\newcommand{\normsub}{\mathbin{\triangleleft}}
\newcommand{\normsub}{\lhd}
\newcommand{\normsubeq}{\trianglelefteq}
\newcommand{\uds}[1]{\mathcal{U}(#1)}
\newcommand{\inv}[1]{#1^{-1}}
\newcommand{\ima}{\text{Im }}
\newcommand{\isom}{\simeq}
\newcommand{\autom}[1]{\text{Aut}(#1)}
\newcommand{\gen}[1]{\langle#1\rangle}
\newcommand{\biy}[1]{\text{Biy}(#1)}
\newcommand{\quo}[2]{{\left.\raisebox{.2em}{$#1$}\middle/\raisebox{-.2em}{$#2$}\right.}}
\newcommand{\appad}[1]{\mathrel{\mkern#1}}
\newcommand{\setdef}[1]{\left\{#1\right\}}
\newcommand{\sdf}[1]{\left\{#1\right\}}
\newcommand{\car}[1]{\mathrm{car}(#1)}
\newcommand{\aut}{\mathrm{Aut}}
\newcommand{\gal}{\mathrm{Gal}}
\newcommand{\actar}{\curvearrowright}
\newcommand{\vabs}[1]{\left|#1\right|}
\newcommand{\rest}[1]{\left.#1\right|}
\renewcommand{\exp}[1]{\mathrm{exp}(#1)}
\newcommand{\mcm}[1]{\mathrm{mcm}\sdf{#1}}
\newcommand{\mtc}[1]{\mathcal{#1}}
\newcommand{\Arg}{\mathrm{Arg}}
%\DeclareMathSymbol{\varprod}{\mathop}{largesymbolsA}{16}

\newcommand{\hr}{\rule{\textwidth}{.4pt}}
\newcommand{\N}{\mathbb{N}}
\newcommand{\Z}{\mathbb{Z}}
\newcommand{\R}{\mathbb{R}}
\newcommand{\Q}{\mathbb{Q}}
\newcommand{\C}{\mathbb{C}}
\newcommand{\A}{\mathbb{A}}
\newcommand{\F}{\mathbb{F}}
\newcommand{\J}{\mathbb{J}}
\newcommand{\B}{\mathbb{B}}
\DeclareMathOperator{\diag}{diag}
\DeclareMathOperator{\traza}{traza}
\renewcommand{\P}{\mathcal{P}}
\renewcommand{\S}{\mathbb{S}}
\newcommand{\ZnZ}{\quo{\Z}{n\Z}}
\newcommand{\ZmZ}{\quo{\Z}{m\Z}}
\newcommand{\0}{\mathbf{0}}
\newcommand{\1}{\mathbf{1}}
\newcommand{\gr}{\text{gr }}
\newcommand\ninf[1]{\left|\left|#1\right|\right|_\infty}
\newcommand\cnj[1]{\overline{#1}}
% \renewcommand\Re{\mathrm{Re}}
% \renewcommand\Im{\mathrm{Im}}
\newcommand\fn[5]{
    \begin{center}
        \begin{tikzpicture}[node distance=0.5cm and 2cm, auto]
            \node (A) [above of =B] {#2};
            \node (B) {#4};
            \node (E) [right= of B] {#5};
            \node (D) [above of =E] {#3};
            \draw[->] (A) to node {#1}(D);
            \draw[|->] (B) to node {}(E);
        \end{tikzpicture}\\
    \end{center}
}

\newcommand{\bfrac}[2]{\left(\frac{#1}{#2}\right)}

%%%% MATRICES
\newcommand\mx[1]{\begin{bmatrix*}#1\end{bmatrix*}}

\newcommand{\suc}[3][1]{{#2}_{#1}, \ldots, {#2}_{#3}}
\newcommand{\sucn}[1]{\suc{#1}{n}}

\renewcommand\qedsymbol{$\diamondsuit$}

\title{Apuntes de Variable Compleja}
\author{Rafael S\'{a}nchez S\'{a}nchez}
\begin{document}
\begin{titlepage}
	\includepdf{portada.pdf}
\end{titlepage}


Revisión del \today $ $ a las \currenttime.

\begin{center}

\end{center}

\tableofcontents
\part{Primer parcial}
% !TeX root = ../compleja.tex

\chapter{N\'umeros complejos y funciones}

\section{Operaciones aritm\'eticas en el cuerpo de los complejos}

En este curso estudiaremos el cuerpo $\C$ formado por el cuerpo de los números complejos.
\begin{dfn}[Número complejo. Parte real e imaginaria]
    Diremos que un número $z$ es \textbf{complejo} cuando sea una tupla de la forma $(a, b)$ (o equivalentemente $(a + bi)$) con $a, b \in \R$. Llamaremos parte real $\Re(z)$ y parte imaginaria $\Im(z)$ a cada escalar $a$ y $b$ de la tupla respectivamente.
\end{dfn}
\begin{dfn}[$\C$]
    Definimos $\C$ como el \textit{cuerpo} conformado por la estructura $\gen{\mtc C, +, \cdot}$, donde $\mtc C = \sdf{(a, b) \mid \forall a,b\ \in \R}$ y las operaciones $+$, $\cdot$ definidas como:
    $$
        (a, b) + (c, d) = (a + c, b + d)
    $$
    $$
        (a, b) \cdot (c, d) = (ac - bd, ad + bc)
    $$
\end{dfn}

\begin{pro}[$\C$ es un cuerpo]
    La estructura $\C = \gen{\mtc C, +, \cdot}$ definida anteriormente satisface las condiciones de ser un cuerpo.
\end{pro}
\begin{proof}
    Se deja al lector.
\end{proof}

\begin{eg}[C\'alculo de un inverso en $\C$]
    Por construcción el neutro de la suma en $\C$ es $(0, 0)$ y el de la multiplicación es $(1, 0)$. Vamos a buscar la expresión del inverso de un complejo $z = (a, b)$. Para ello buscamos resolver el sistema:
    $$
        (a, b) \cdot (x, y) = (1, 0)
    $$
    es decir:
    \begin{align*}
        ax - by &= 1\\
        -bx + ay &= 0
    \end{align*}
    que tiene solución cuando $a^2 + b^2 \neq 0$. Finalmente obtenemos:
    $$
        (x, y) = \left( \frac{a}{a^2+b^2}, \frac{-b}{a^2+b^2} \right)
    $$
    que como veremos más adelante implica que:
    $$
        z^{-1} = \frac{\cnj z}{\vabs{z}^2}
    $$
\end{eg}

Además, tiene sentido que si construimos $\C$ a partir de tuplas $\R\times \R$ entonces $\R \subseteq \C$. Nos interesa definir exactamente como, para ello establecemos la \textit{función de inclusión:} $\iota$\\
\fn{$\iota$}{$\R$}{$\C$}{$a$}{$(a, 0)$}

\begin{obs}
    En ocasiones usaremos indistintamente $(a, 0) \equiv a$ cuando un número complejo solo tenga parte real.
\end{obs}

Veremos más adelante que también podemos calcular las raíces de números complejos. En particular para las raíces cuadradas se reduce a resolver el sistema que se deduce de $(x, y)^2 = (a, b)$, es decir:
\begin{align*}
    x^2 - y^2 &= a\\
    2xy &= b
\end{align*}
que es un sistema resoluble, sin embargo la expresión no es nada agradable. Usaremos la \textit{representación polar} de los números complejos para simplificar esta tarea.

\subsection{Conjugaci\'on}
    \begin{dfn}[Conjugado de un número complejo]
        Definimos el \textbf{conjugado} de un número complejo $z = (a, b) \in \C$ como $\cnj z = (a, -b)$.
    \end{dfn}

    \begin{dfn}[Módulo de un número complejo]
        Definimos el \textbf{módulo} de un número complejo $z = (a, b) \in \C$ como $\rho = \vabs{z} = \sqrt{a^2 + b^2} \in \R$.
    \end{dfn}

    \begin{pro}[Propiedades del conjugado de un número complejo]
        Sea $z = (a, b) \in \C$:\\
        \begin{enumerate}[(1)]
            \item $z \cdot \cnj z = (a^2 + b^2) = \vabs{z}^2$
            \item $\cnj{z \cdot w} = \cnj z \cdot \cnj w$
            \item $\cnj{z + w} = \cnj z + \cnj w$
            \item $\cnj{\cnj z} = z$
            \item $z + \cnj z = 2 \cdot \Re(z)$
            \item $z - \cnj z = 2 \cdot \Im(z)$
            \item $z = \cnj z \iff z \in \R$
            \item Si $w \neq 0$, $\cnj{\frac{z}{w}} = \frac{\cnj z}{\cnj w}$
        \end{enumerate}
    \end{pro}
    \begin{proof} Consideraremos $z = (a, b)$ y $w = (c, d)$.
        \begin{enumerate}
            \item $$(a, b) \cdot (a, -b) = (a \cdot a - b \cdot(-b), a \cdot (-b) + b \cdot a) = (a^2 + b^2, 0)$$
            \item
            \begin{align*}
                \cnj{z\cdot w} &= \cnj{(ac - bd, ad + bc)} = (ac - bd, -ad -bc) =\\
                &= (ac - (-b) (-d), a(-d) + c(-b)) = (a, -b) \cdot (c, -d) = \cnj z \cdot \cnj w
            \end{align*}
            \item $$
                \cnj{z + w} = \cnj{(a + c, b + d)} = (a + c, -b -d) = (a, -b) + (c, -d) = \cnj z + \cnj w
                $$
            \item $$\cnj{\cnj{z}} = \cnj{(a, -b)} = (a, b) = z$$
            \item $$z + \cnj z = (a, b) + (a, - b) = (2a, 0) = 2 \Re(z)$$
            \item $$z - \cnj z = (a, b) - (a, - b) = (0, 2b) = 2 \Im(z)$$
            \item $$ z = \cnj z \iff a = a \text{ y } b = -b \iff b = 0 \iff z \in \R$$
            \item Vamos a demostrar primero que $\cnj{\frac{1}{w}} = \frac{1}{\cnj w}$:
            $$
                1 = w \cdot \frac{1}{w} \implies 1 = \cnj 1 = \cnj{w \cdot \frac{1}{w}} = \cnj{w} \cdot \cnj{\frac{1}{w}}
            $$
            y como $\C$ es un cuerpo sabemos que $\cnj{w}^{-1}$ existe, por tanto, multiplicando a ambos lados por $\cnj{w}^{-1}$ obtenemos:
            $$
                \frac{1}{\cnj{w}} = \cnj{\frac{1}{w}}
            $$
            Con este resultado ya es directo demostrar que:
            $$
                \cnj{\frac{z}{w}} = \cnj z \cdot \cnj{\frac{1}{w}} = \cnj z \frac{1}{\cnj w} = \frac{\cnj z}{\cnj w}
            $$
        \end{enumerate}
    \end{proof}
\subsection{Desigualdad triangular}
    En esta sección vamos a ver algunas propiedades de los números complejos así como la desigualdad triangular y su generalización.
    \begin{pro}[Propiedades del módulo de un número complejo]
        De nuevo consideramos $z = (a, b)$, $w = (c, d)\in \C$, $\sdf{z_i=(a_i, b_i)}_{i\in \N} \subset \C$.
        \begin{enumerate}
            \item $\vabs{z \cdot w} = \vabs{z} \cdot \vabs{w}$.
            \item $\vabs{z} = \vabs{-z} = \vabs{\cnj z} = \vabs{-\cnj z}$.
            \item $\vabs{\Re(z)} \leq \vabs{z}$ y $\vabs{\Im(z)} \leq \vabs{z}$.
            \item $\vabs{z + w} \leq \vabs{z} + \vabs{w}$. (Desigualdad triangular).
            \item $\vabs{z} - \vabs{w} \leq \vabs{z + w}$.
            \item $\vabs{ \vabs{z} - \vabs{w} } \leq \vabs{z + w}$.
            \item $\vabs{z_1 + \cdots + z_n} \leq \vabs{z_1} + \cdots + \vabs{z_n}$. (Desigualdad triangular generalizada).
        \end{enumerate}
    \end{pro}
    \begin{proof} En la gran mayoría de apartados se procede a demostrar la relación de los cuadrados, ya que al ser $f(x) = \sqrt{x}$ una función estrictamente creciente para números positivos, basta tomar raíces a ambos lados y llegamos a relación original.
        \begin{enumerate}
            \item $$\vabs{z \cdot w}^2 = (ac-bd)^2+(ad+bc)^2 = a^2c^2 \cancel{-2abcd} + b^2d^2 + a^2d^2 +\cancel{2abcd} + b^2c^2 = (a^2+b^2)(c^2+d^2) = \vabs{z}^2 + \vabs{w}^2$$
            \item Los cuatro términos se pueden abreviar en $\tilde z = (\pm a, \pm b)$ y entonces: $$\vabs{\tilde z} = \sqrt{(\pm a)^2 + (\pm b)^2} = \sqrt{a^2 + b^2} = \vabs{z}$$
            \item Como $\vabs{z}^2 = \vabs{\Re(z)}^2 + \vabs{\Im(z)}^2$ entonces es claro que:
            $$
                \vabs{\Re(z)}^2 \leq \vabs{\Re(z)}^2 + \vabs{\Im(z)}^2 \text{ y } \vabs{\Im(z)}^2 \leq \vabs{\Re(z)}^2 + \vabs{\Im(z)}^2
            $$
            \item
            \begin{align*}
                \vabs{z + w} &\leq \vabs{z} + \vabs{w}\\
                \vabs{z + w}^2 &\leq (\vabs{z} + \vabs{w})^2\\
                (z+w)\cnj{(z+w)} &\leq \vabs{z^2} + \vabs{w^2} + 2\vabs{z}\vabs{w}\\
                (z+w)(\cnj z + \cnj w) &\leq \vabs{z^2} + \vabs{w^2} + 2\vabs{z}\vabs{w}\\
                \vabs{z^2} + \vabs{w^2} + z\cnj w + \cnj z w &\leq \vabs{z^2} + \vabs{w^2} + 2\vabs{z}\vabs{w}\\
                \cancel 2\cdot\Re(z\cnj w) = z \cnj w + w \cnj z &\leq \cancel 2 \vabs{z} \vabs{w}\\
                \Re(z\cnj w) &\leq \vabs{z}\vabs{w} = \vabs{z}\vabs{\cnj w} = \vabs{z\cnj w}
            \end{align*}
            \item
            $$
                \vabs{z} = \vabs{z + w - w} \leq \vabs{z + w} + \vabs{-w} = \vabs{z + w} + \vabs{w} \implies \vabs{z} - \vabs{w} \leq \vabs{z + w}
            $$
            Análogamente demostramos la segunda desigualdad expresando $\vabs{w} = \vabs{w + z - z}$.
            \item Directa de la propiedad anterior y la desigualdad triangular.
            \item Por inducción sobre el número de términos del sumatorio de la desigualdad triangular. Sabemos que el caso base se cumple ($n=2$), ahora suponemos que se cumple para $n$ sumandos y lo probamos cierto para $n+1$, es decir, suponemos que es cierto:
            $$
                \vabs{\sum_{i=1}^{i=n} z_i} \leq \sum_{i=1}^{i=n} \vabs{z_i}
            $$ y entonces:
            $$
                \vabs{\sum_{i=1}^{i=n+1} z_i} = \vabs{z_{n+1} + \sum_{i=1}^{i=n} z_i} \leq \vabs{z_{n+1}} + \vabs{\sum_{i=1}^{i=n} z_i} \leq \vabs{z_{n+1}} + \sum_{i=1}^{i=n} \vabs{z_i} = \sum_{i=1}^{i=n+1} \vabs{z_i}
            $$
        \end{enumerate}
    \end{proof}
\subsection{Representaci\'on polar}
    Los números complejos, al estar construidos como un par ordenado de números reales se pueden representar geométricamente en el plano complejo (similar a $\R^2$). De esta forma, el número $z = a + bi = (a, b)$ se puede representar como el punto $(a, b)$ de $\R^2$.\\\\
    Además, toman especial importancia los puntos de la circunferencia unidad, (los que satisfacen la ecuación $x^2 + y^2 = 1$). Usando un poco de trigonometría podemos obtener fácilmente que los puntos del plano complejo que se corresponden con los de la circunferencia unidad son $(\cos \theta, \sin \theta)$.

    \begin{obs}
        Los números complejos no tienen una expresión única en esta forma ya que $(\cos \theta, \sin \theta) = (\cos (\theta + 2k\pi), \sin  (\theta + 2k\pi) )$ con $k \in \Z$.\\
        El $(0, 0)$ no tendría sentido expresarlo de forma polar, lo denominaremos simplemente $0$.
    \end{obs}

    \begin{center}
        \begin{tikzpicture}
            \path (63: 1.5) coordinate (P);
            \path (120:1) coordinate (Q);
            \path (0: 0.3) coordinate (R);
            \path (0: 0.5) coordinate (R1);
            \draw (0, 0) ellipse (1cm and 1cm);
            \draw (-2, 0) -- (2, 0);
            \draw (0, -1.2) -- (0, 1.2);
            \filldraw[black] (P) circle (1pt) node[anchor=west] {$z$};
            \filldraw[black] (Q) circle (1pt) node[anchor=east] {$(x, y)$};
            \draw (0, 0) -- (P);
            \draw[->] (R) arc (0:63:0.3) node[anchor=west] {$\theta$};
            \draw[->] (R1) arc (360:63:0.5) node [anchor=south] {$\sigma$};
        \end{tikzpicture}
    \end{center}

    De esta forma, cualquier número complejo $z$ podrá expresarse como $z = \vabs{z}(\cos \theta, \sin \theta)$. Usualmente denominaremos a $\theta$ como el argumento de $z$, es decir, $\theta=\arg(z)$. Además, es común encontrar el módulo expresado con la letra $\rho$. Con esto podremos expresar $z$ como:\\
    $$
        z = \rho(\cos \theta, \sin \theta) = \rho(\cos \theta + i \sin \theta)
    $$

    Cuando definamos la exponencial compleja veremos que podemos expresar $z$ de la forma $z = \rho e^{i\theta}$.

    \begin{obs} %%TODO: Referenciar la imagen
        Como se aprecia en la figura superior, el argumento de $z$ no es único, tanto $\sigma$ como $\theta$ son el argumento de z (debido a que $\sigma = \theta - 2\pi$).\\

        Debido a que $\arg(z)$ es una función multievaluada (si $\arg(z) = \theta \implies \arg(z) = \sdf{\theta + 2n \pi \mid n \in \Z}$), en ocasiones querremos usar una función que escoja un solo argumento de ellos. La denominamos \textbf{argumento principal} ($\Arg{z}$), y la definimos restringiendo la imagen de $\arg{z}$ al intervalo $\left[0, 2\pi\right)$ o $\left[-\pi, pi\right)$.\\

        Sin embargo, aun habiendo determinado un argumento principal, en ninguno de los dos casos la función $\Arg()$ es continua en todo punto. Se puede ver que no existe una determinación continua del argumento.
    \end{obs}

    Esta forma de representar los números complejos nos simplificará varias tareas, como multiplicarlos o hallar raíces.
%%%%%%%%%%%%%%%%%%%%%%%%%%%%%%%%%%%%%%%%%%%%%%%%%%%%%%%%%%%%%%%%%%%%%%%%%% 30/01

    \begin{obs}
        Sea $\rho(\cos \theta, \sin \theta) = z = a + bi \implies \frac{b}{a} = \tan{\theta}$. Esto no significa que $\theta = \tan^{-1}\left(\frac{b}{a}\right)$ ya que la imagen de $\tan^{-1} \subseteq \left(-\frac{\pi}{2}, \frac{\pi}{2}\right)$.
    \end{obs}

    En general, encontrar el argumento de un número complejo no es tarea fácil, sin embargo en ciertos casos puede comprobarse a mano.

    \begin{eg}[Hallar la expresión polar de un número complejo. Hallar su inverso]
        Sea $z = 1-i \in \C$, vamos a hallar su expresión polar.
        $$
            \vabs{z} = \sqrt{1^2 + (-1)^2} = \sqrt 2
        $$
        Por tanto, $z = \sqrt{2} \cdot \left( \frac{1}{\sqrt 2} - \frac{i}{\sqrt 2} \right) = \sqrt 2 (\cos \theta + i \sin \theta)$. Resolviendo el sistema:
        \begin{align*}
            \cos \theta &= \frac{1}{\sqrt 2}\\
            \sin \theta &= - \frac{1}{\sqrt 2}
        \end{align*}
        Y al resolverlo hallamos $\theta = -\frac{\pi}{4}$ y entonces:
        $$
            z = \sqrt{2} (\cos \left(-\frac{\pi}{4}\right) + i \sin \left(-\frac{\pi}{4}\right))
        $$

        Si ahora queremos hallar su inverso, vamos a usar la identidad $z^{-1} = \frac{\cnj z}{\vabs{z}^2}$. Por tanto:
        $$
            z^{-1} = \frac{1}{2} (1 + i)
        $$
    \end{eg}

\subsection{Ra\'ices y potencias}

    Ya vimos al comienzo del capítulo que la forma de hallar la raíz cuadrada de un número complejo en la primera expresión que dimos del mismo era poco agradable. Podemos ver que la forma polar nos simplifica el cálculo.\\
    Para ello vamos a comenzar observando como se comporta la expresión polar frente a la multiplicación de complejos.\\

    Consideremos $z = \rho(\cos \theta, \sin \theta)$ y $w = \rho'(\cos \theta', \sin \theta')$. Entonces:
    $$
        z \cdot w = \rho \cdot \rho' (\cos \theta \cos \theta' - \sin \theta \sin \theta', \cos\theta \sin \theta' + \sin \theta \cos \theta') = \rho \cdot \rho' (\cos (\theta + \theta') + i\sin(\theta+\theta'))
    $$
    Es decir, cuando multiplicamos dos complejos en forma polar, el resultado es otro número complejo cuyo módulo es el producto de los módulos y el argumento es la suma de los argumentos. Esta propiedad simplifica el cálculo de raíces y potencias.

    \begin{eg}[Cálculo de una potencia de un número complejo]
        Sea $z = (\cos \theta, \sin \theta)$, entonces, usando lo que acabamos de ver:
        $$
            z^n = (\cos \theta, \sin \theta)^n = \cos(n\theta) + i \sin(n\theta)
        $$
    \end{eg}

    \begin{eg}[Cálculo de las raíces de un número complejo]
        Sean $z, w \neq 0 \in \C$, queremos resolver la ecuación $z^n = w$, es decir, encontrar las raíces n-ésimas de $w$.\\
        Podemos suponer que $w = \rho(\cos \theta + i \sin \theta)$ y $z = r (\cos t + i \sin t)$. Entonces queremos que:
        $$
            z^n = w \implies r^n (\cos(nt) + i \sin(nt)) =  \rho(\cos \theta + i \sin \theta)
        $$
        y entonces:
        \begin{align*}
            r &= \sqrt[n]{\rho}\\
            \theta + 2\pi k &= nt
        \end{align*}
        Simplificando la última identidad, obtenemos:
        $$
            t = \frac{\theta}{n} + 2\pi \frac{k}{n},\ k=0, \ldots, n-1
        $$
        y por tanto resultan $n$ raíces. (Para $k=n$ el complejo es el mismo que para $k = 0$).
    \end{eg}
% \section{Topolg\'ia del plano complejo}
% \section{Esfera de Riemann}
% \section{Funciones complejas}
% \section{L\'imites y continuidad}

% % !TeX root = ../galois.tex

\chapter{Extensiones de cuerpos}

\section{Grados de cuerpos}

\begin{dfn}[Extensión]
    Sean $K, E$ cuerpos, decimos que $E$ es una \textbf{extensión} de $K$ (denotado por $E/K$) si $K$ es un subcuerpo de $E$.
\end{dfn}

\begin{eg}[Extensiones]$ $
    \begin{itemize}
        \item $\C/\Q$
        \item $\R/\Q$
        \item $\C/\R$
        \item $\Q[\sqrt{n}]/\Q$ con $n \neq$ de un cuadrado perfecto.
    \end{itemize}
\end{eg}

\begin{pro}[Extensión como espacio vectorial]
    Si $E$ es una extensión de $K$, entonces $E$ es un espacio vectorial sobre $K$.
\end{pro}
\begin{proof}
    Basta interpretar el producto por escalares $\cdot: K \times E \to E$ como la restricción del producto sobre $E \times E$ a $K$. La suma está bien definida por ser $E$ un grupo abeliano con la suma.
\end{proof}

\begin{dfn}[Grado de una extensión]
    Sea $E/K$ una extensión, el grado de la extensión es $|E:K| = \dim_K E$, que coincide con la dimensión del espacio vectorial que define $E$ sobre $K$.
\end{dfn}

\begin{dfn}[Extensión finita]
    Sea $E/K$ una extensión, es \textbf{finita} si y sólo si $\exists \setdef{a_1, \ldots, a_n} \subseteq E$ tales que forman una $K$-base. Es equivalente a decir que $\dim_K E = n < \infty$.
\end{dfn}

\begin{lm}[Extensión de grado 1]
    Sea $E/K$ una extensión:
    $$
        |E:K| = 1 \iff E = K
    $$
\end{lm}

\begin{proof}$ $
    \begin{itemize}
        \item[$\implies$] Si $|E:K| = \1$, entonces $\exists e \in E$ tal que $\setdef{e}$ es una $K$-base. Por tanto: $\1 = k \cdot e$ con $k \in K \implies e = k^{-1} \implies e \in K \implies E = K$.
        \item[$\ \Longleftarrow \ $] $\setdef{\1}$ es $K$-base de $K$. $\dim_K K = 1 = |K:K|$.
    \end{itemize}
\end{proof}

\begin{thm}[Transitividad de grados]\label{thm:transgr}\label{thm:2.1}% 2.1
    Sea una extensión $E/K$ y un subcuerpo $L$ intermedio $K \subseteq L \subseteq E$, entonces la extensión $E/K$ es finita si y sólo si $|E:L| < \infty,\ \land |L:K| < \infty$, y en tal caso: $|E:K| = |E:L|\cdot |L:K|$.
\end{thm}

\begin{proof}
    Supongamos $\dim_K E = r < \infty$, y $\setdef{e_1, \ldots, e_r}$ una $K$-base de $E$, entonces $\setdef{e_1, \ldots, e_r}$ es un $L$-sistema generador de $E \implies E/L$ es finita.\\
    Como $K \subseteq L \subseteq E$, $L$ es un $K$-subespacio vectorial de $E$, en particular $\dim_K L \leq dim_K E < \infty$. Si $E/L$ y $L/K$ son finitas, cogemos $\setdef{b_1, \ldots, b_m}$ una $L$-base de $E$ y $\setdef{a_1, \ldots, a_n}$ una $K$-base de $E$.\\
    Queremos ver que $\setdef{a_ib_j \mid 1 \leq i \leq n,\ 1 \leq j \leq m}$ es una $K$-base de $E$ (en particular con esto habremos probado que $|E:K| = |E:L|\cdot |L:K|$). Sabemos que:
    $$
        x \in E, x = \sum_{j=1}^m l_j b_j,\ l_j \in L
    $$
    pero además
    $$
        l_j = \sum_{i = 1}^{n} k_{ij} a_i,\ k_{ij} \in K \implies x = \sum_{1 \leq i \leq n,\ 1 \leq j \leq m} k_{ij}a_ib_j,\ k_{ij} \in K
    $$
    Faltaría ver que $\setdef{a_ib_j}$ es $K$-libre.
    $$
        \sum_{1 \leq n,\ 1 \leq i\leq j \leq m} k_{ij}a_ib_j = 0 \implies \sum_j l_j b_j = 0 \implies l_j = 0 \implies \sum_i k_{ij}a_i = 0 \implies k_{ij} = 0 \forall i \forall j
    $$
\end{proof}

\begin{dfn}[Menor subanillo y subcuerpo]Sea $E/K$ una extensión y sea $a \in E$.
    \begin{itemize}
        \item Denotaremos por $K[a]$ al \textbf{menor subanillo} de $E$ que contiene a $K$ y a $a \in E$. Se puede probar que $K[a] = \setdef{f(a) \forall f \in K[x]}$.
        \item Denotaremos por $K(a)$ al \textbf{menor subcuerpo} de $E$ que contiene a $K$ y a $a \in E$. Se puede probar que $K(a) = \setdef{\frac{f(a)}{g(a)} \forall f,g \in K[x] \mid g(a) \neq 0}$.
    \end{itemize}
\end{dfn}

\begin{eg}$ $
    \begin{itemize}
        \item $\Q[\sqrt{n}] = \setdef{a + \sqrt{n} \mid a, b\in \Q}$. En este caso $\Q[\sqrt{n}] = \Q(n)$.
        \item $X \subseteq E$, $K(X)$ es el menor subcuerpo de $E$ que contiene a $K$ y a $X$. $K(X)$ se obtiene al \textit{adjuntar} $X$ a $K$.
    \end{itemize}
\end{eg}

\begin{obs}
    En general $K[a] \subseteq K(a)$, pero hay en casos en los que la igualdad no se cumple.
\end{obs}

\begin{dfn}[Extensión simple]
    Sea $E/K$ una extensión, es \textbf{simple} si $\exists a \in E$ tal que $E = K(a)$.
\end{dfn}

\begin{eg}[Extensión simple]$ $
    \begin{itemize}
        \item $\C/\R$, ya que $\C = \R(i)$.
        \item Con $p\neq q$ primos, $\Q(\sqrt{p}, \sqrt{q})$ es simple, ya que se puede demostrar que $\Q(\sqrt{p}, \sqrt{q}) = \Q(\sqrt{p} + \sqrt{q})$.
    \end{itemize}
\end{eg}

\begin{pro}[Dimensión de un cuerpo finito]
    Sea $K$ un cuerpo finito, entonces $|K| = p^n$ con $p$ primo.
\end{pro}

\begin{proof}
    Sea $K$ un cuerpo y $F$ su cuerpo primo, sabemos que $F$ es isomorfo a algún $\F_p$ con $p$ un primo. Además, $K/F$ es una extensión. Y como $K$ es subespacio vectorial $|K:F| = \dim_F K = n$. Entonces $K \isom F^n \implies |K| = |F|^n = p^n$.
\end{proof}

\section{Extensiones algebraicas y trascendentes}

\begin{dfn}[Extensión algebraica. Extensión trascendente]
    Sea $E/K$ una extensión.
    \begin{itemize}
        \item Sea $a \in E$, $a$ es \textit{algebraico} si $\exists f(x) \neq 0 \in K[x]:\ f(a) = 0$. $E/K$ es una \textbf{extensión algebraica} si todo $a \in E$ es algebraico sobre $E$.
        \item Sea $a \in E$, $a$ es \textit{trascendente} si no es algebraico. $E/K$ es una \textbf{extensión trascendente} si existe $a \in E$ trascendente sobre $E$.
    \end{itemize}
\end{dfn}

\begin{eg}[Extensiones algebraicas y trascendentes]$ $
    \begin{itemize}
        \item $K/K$ es algebraica. Todo elemento de $K$ es raíz de $x - k \in K[x]$.
        \item $\Q(\sqrt{n})/\Q$ es algebraica.
        \item $\R/\Q$ es trascendente. $e$ y $\pi$ son trascendentes.
        \item Sea $K[t]$ un dominio de integridad- Podemos construir su cuerpo de fracciones:
        $$
            K(t) = \setdef{\frac{f(t)}{g(t)} \mid f,g \in K[t],\ g(t) \neq 0}
        $$
        Entonces $K(t)/K$ es trascendente. ($t$ siempre es trascendente).
    \end{itemize}
\end{eg}

\begin{pro}[Extensiones y cuerpos intermedios]
    Sea $E/K$ una extensión y $K \subseteq L\subseteq E$ un cuerpo intermedio:
    \begin{enumerate}
        \item $E/K$ es algebraica $\iff L/K$ y $E/L$ son algebraicas.
        \item Si $E/L$ es trascendente, entonces $E/K$ es trascendente.
    \end{enumerate}
\end{pro}

\begin{proof}
    Se deja como ejercicio.
\end{proof}

\begin{thm}[Extensiones finitas y algebraicas]\label{thm:finitaalgebraica}\label{thm:2.2} %2.2
    Toda extensión finita es algebraica.
\end{thm}
\begin{proof}
    Sea $E/K$ una extensión, $a \in E$, queremos ver que es raíz de $0 \neq f(x) \in K[x]$. Suponemos $|E : K| = n$ con un $K$-sistema: $\setdef{1, a_1, \ldots, a^{n-1}} \subseteq E$ con $n$ elementos. Entonces, el sistema puede ser:
    \begin{itemize}
        \item [$K$-ligado] Existen $k_i \in K$ no todos nulos tales que $k_0 + k_1a + \ldots + k_{n-1}a^{n-1} = 0$, entonces $a$ es raíz de $f(x) = k_0 + \ldots + k_{n-1}x^{n-1}$.
        \item [$K$-libre] Como $\dim_K E = n$, $\setdef{1, a_1, \ldots, a^{n-1}, a^n}$ es $K$-ligado, y de nuevo, $a$ es raíz de $f(x) = k_0 + \ldots + k_{n-1}x^{n-1} + k_n x^n$.
    \end{itemize}
\end{proof}

\section{Teorema del elemento algebraico}

\begin{thm}[Teorema del elemento algebraico]\label{thm:elemalg}\label{thm:2.3} %2.3
    Sea $E/K$ una extensión, $a \in E$ un elemento algebraico sobre $K$.
    \begin{enumerate}
        \item Existe un único polinomio irreducible mónico $p \in K[x]$ tal que $p(a) = 0$.
        \item Si $q \in K[x]$ y $q(a) = 0$, entonces $p \divides q$.
        \item $K(a) = \setdef{f(a) \mid f \in K[x]} = K[a]$.
        \item Si $\delta(p) = n$, entonces $\setdef{1, a, \ldots, a^{n-1}}$ es una $K$-base de $K(a)$. En particular, $|K(a) : K| = \delta(p)$ y $K(a) = \setdef{k_0 + k_1 a +\ldots + k_{n-1} a^{n-1} \mid k_i \in K}$.
    \end{enumerate}
\end{thm}

\begin{proof}
    ($\cdots$) (Es muy tarde ahora mismo para pasar esto a limpio, que alguien me mate por favor.) % TODO
\end{proof}

\begin{dfn}[Polinomio mínimo]
    Sea $E/K$ una extensión y $a\in E$ algebraico sobre $K$, al único polinomio mónico e irreducible $p \in K[x]$ dado por el teorema \ref{thm:elemalg} se le llama \textbf{polinomio mínimo} o \textbf{polinomio irreducible} de $a$ sobre $K$ y escribimos $p = Irr(K, a)$.
\end{dfn}

\begin{obs}
    Sea $b = k_0 + \ldots + k_{n-1} a^{n-1} \in K(a)$, ¿cómo se expresa $b^{-1}$ en la misma base?\\
    Consideramos $f(x) = k_0 + \ldots + k_{n-1}x^{n-1} \in K[x]$, con $f(a) = b \neq 0$ y $\delta(f) \leq \delta(p)$, siendo $p$ el polinomio irreducible. Entonces $mcd(f, g) = 1$. Por la identidad de Bezout:
    $$
        \exists h, g \in K[x]:\ 1 = fh + gp \implies \text{ (evaluando en $a$)} 1 = f(a)h(a) + g(a)p(a) \implies 1 = f(a)h(a) = b h(a) \implies b^{-1} = h(a)
    $$
\end{obs}

\begin{eg}[Ejercicio tipo]
    ($\cdots$) (Es aún más tarde. \url{https://www.youtube.com/watch?v=I_6Gej1m4SU}) % TODO
\end{eg}

\begin{thm}[Extensión por varios elementos algebraicos]\label{thm:extalg}\label{thm:2.4} %2.4
    Sea $E/K$ una extensión, y sea $\mbf{a} = (a_1, \ldots, a_n)$ con $a_i \in E$ algebraicos sobre $K$, entonces $K(\mbf{a})/K$ es finita. En particular, $K(\mbf{a})/K$ es algebraica.
\end{thm}

\begin{proof}
    Por inducción sobre $n$. Si $n=1$ por el teorema del elemento algebraico (teorema \ref{thm:elemalg}) sabemos que $|K(a_1)/K| = \delta(Irr(K, a_1)) < \infty$.\\
    Veamos el caso $n>1$. Sea $L = K(a_1, \ldots, a_{n-1})$, entonces $K(\mbf{a}) = L(a_n)$. Por la hipótesis de inducción $L/K$ es finita, por el teorema \ref{thm:elemalg} $L(a_n)/L$ es finita y por tanto, por el teorema de transitividad de grados (teorema \ref{thm:transgr}) $L(a_n)/K$ es finita, donde $L(a_n)$ era $K(\mbf{a})$. La segunda parte se sigue directamente aplicando el teorema \ref{thm:finitaalgebraica}.
\end{proof}

\begin{ex}[H2.7]
    Dada $E/K$ una extensión, prueba que $L = \setdef{e \in E \mid e \text{ es algebraico sobre } K} \supseteq K $ es un cuerpo. Sea $\A \subset \C$ los elementos algebraicos sobre $\Q$, prueba que $\A/\Q$ es una extensión de grado infinito\\
    Sean $a, b \in L$ cualesquiera, entonces $a, b \in K(a, b)^\times$ y por el teorema \ref{thm:extalg} $K(a,b)/K$ es algebraica. Como $a \pm b$ y $ab^{\pm 1} \in K(a, b)$, entonces son algebraicos sobre $K \implies (a \pm b), (ab^{\pm 1}) \in L$, con lo que es cerrado por ambas operaciones y $L$ es un cuerpo.\\
    Por la primera parte del ejercicio, sabemos que $\Q \subseteq \A \subseteq \C$ es un subcuerpo intermedio, de $\C/\Q$, y por definición $\A/\Q$ es algebraica.\\
    Por el criterio de Einsestein (teorema \ref{thm:Einsestein}), para cada $n \in \N^\times$, $x^n - 2$ es irreducible en $\Q[x]$. $\sqrt[\leftroot{-3}\uproot{3}n]{2}$ es solución por tanto es algebraico y por el teorema \ref{thm:elemalg} $|\Q(\sqrt[\leftroot{-3}\uproot{3}n]{2})/\Q|=n$.\\
    Como podemos hacerlo $\forall n \in \N^\times$, hemos comprobado que $\A/\Q$ es no finita.
\end{ex}

\begin{ex}[H2.6]
    ($\cdots$) (Yep. \url{https://www.youtube.com/watch?v=pwSsT8IU0WE})
\end{ex}

\section{Isomorfismos de cuerpos}

Ya definimos qué es un homomorfismo de cuerpos en la sección \ref{section:homomorfismos}. En esta sección vamos a ampliar los conocimientos cuando aplicamos los homomorfismos específicamente a cuerpos.

\begin{obs}
    Si $\car{E} \neq \car{K}$ entonces: $Hom(E, K) = \varnothing = Hom(K, E)$, con $Hom(X, Y) = \setdef{\varphi: X \to Y \mid \varphi \text{ es un homomorfismo de cuerpos}}$. Además, si $K$ es finito, $End(K) = Aut(K)$ (conjunto de endomorfismos y automorfismos respectivamente), pero en general: $End(K) \subsetneq Aut(K)$. %TODO: ¿ esta bien la dirección ?
\end{obs}

\begin{obs}
    Sea $\varphi \in End(K)$ y $F$ es el cuerpo primo de $K$, entonces $\varphi(a) = a,\ \forall a \in F$. De forma más general, si $\varphi \in Hom(E, K)$ y $E, K$ tienen el mismo cuerpo primo, $\varphi(a) = a$, por ejemplo:
    $Aut(\F_p) = \setdef{id}$, $Aut(\Q) = \setdef{id}$.
\end{obs}

En ocasiones querremos saber como extender un isomorfismo de cuerpos a una extensión de dichos cuerpos. Vamos a ver un lema y un teorema.

\begin{lm}[Restricción de un isomorfismo de cuerpos]
    Sea $E_1/K_1$ una extensión, $\theta: E_1 \to E_2$ un isomorfismo de cuerpos, y sea $K_2 = \theta(K_1)$, entonces:
    \begin{enumerate}
        \item $E_2/K_2$ es una extensión y $|E_1:K_1| = |E_2:K_2|$.
        \item Sean $a_1, \ldots, a_n \in E_1$. $\theta(K_1(a_1, \ldots, a_n)) = K_2(\theta(a_1), \ldots, \theta(a_n))$.
        \item $\theta$ se extiende a un isomorfismo de anillos $\theta: K_1[x] \to K_2[x]$, aplicando $\theta$ individualmente a cada coeficiente del polinomio.
    \end{enumerate}
\end{lm}

\begin{proof}
    Como ejercicio.
\end{proof}

\begin{thm}[Extensión de un isomorfismo de cuerpos]\label{thm:2.5} % 2.5
    Sean $E_1/K_1$, $E_2/K_2$ extensiones, $\sigma: K_1 \to K_2$ un isomorfismo de cuerpos, $p_1$ irreducible en $K_1$, $p_2$ irreducible en $K_2$, y $a_1$ y $a_2$ raíces de los polinomios respectivamente, entonces:
    $$
        \sigma \text{ se extiende a } \theta \iff \left.\theta\right|_{K_1} = \sigma
    $$
    Donde $\theta: K_1(a_1) \to K_2(a_2)$ tal que $\theta(a_1) = a_2$.
\end{thm}

\begin{proof}
    A completar.
\end{proof}

\begin{cor} % 2.6, 2.7
    Sea $E/K$ una extensión, $p \in K[x]$ irreducible, entonces:\\

    $$a, b \in E \text{ son raíces de } p \iff \exists \text{ un isomorfismo } \theta: K(a) \to K(b) \text{ tal que } \theta(a) = b \text{ y } \theta(k) = k\ \forall k \in K$$
\end{cor}

\begin{ex}[H2.4 (parte)]
    Halla el grado y base de las siguientes extensiones de cuerpos:
    \begin{enumerate}
        \item $\Q(\sqrt{2}, \sqrt{3}, i)/\Q$.\\
        Consideramos $L = \Q(\sqrt{2}, \sqrt{3}) = \Q(\sqrt{2} + \sqrt{3})$ cuerpo intermedio, y además, $\Q(\sqrt{2}, \sqrt{3}, i) = L(i)$. $i$ es raíz del polinomio $x^2 + 1$ qque no tiene raíces en $L$ y por tanto, $x^2 + 1 = Irr(L, i)$. Por el teorema \ref{thm:elemalg}, $|\Q(\sqrt{2}, \sqrt{3}, i):\Q| = 2 = \delta(Irr(L, i))$.\\

        Además, $|L:\Q| = 4$ por el ejercicio 1 de la hoja 2. Por el teorema \ref{thm:transgr}, $|\Q(\sqrt{2}, \sqrt{3}, i) : \Q| = |L(i):L|\cdot|L:\Q| = 8$. Una vez sabemos el grado, podemos encontrar una $\Q$-base:
        $$
            \setdef{1, \sqrt{2}, \sqrt{3}, \sqrt{6}, \sqrt{2} i, \sqrt{3} i, \sqrt{6} i}
        $$
        $$
            \setdef{1, \alpha, \alpha^2, \alpha^3, i, \alpha i, \alpha^2 i, \alpha^3 i}
        $$

        \item $\Q(\sqrt[4]{2})/\Q(\sqrt{2})$.\\
        $\alpha = \sqrt[4]{2}$ es raíz de $x^2 - \sqrt{2} \in \Q(\sqrt{2})[x]$, y es irreducible porque no tiene raíces en $Q(\sqrt{2})$, (se prueba por reducción al absurdo).\\
        Por tanto: $|\Q(\sqrt[4]{2}):\Q(\sqrt{2})| = 2$ y una base:
        $$
            \setdef{1, \sqrt[4]{2}}
        $$

        \item $\Q(\sqrt{1+\sqrt 3})/\Q$.\\
        Consideramos el cuerpo intermedio $L = \Q(\sqrt{3})$. Es fácil ver que $|L/\Q| = 2$. Falta encontrar el grado de $|\Q(\sqrt{1+\sqrt 3})/L|$. Sea $\alpha = \sqrt{1 + \sqrt 3}$, $\alpha$ es raíz de $x^2 - (1+\sqrt 3)$, que se puede demostrar que es irreducible en $L$. Por tanto, por el teorema \ref{thm:transgr}, $|\Q(\sqrt{1+\sqrt 3}):\Q| = 4$, y la base:
        $$
            \setdef{1, \sqrt 3, \sqrt{1 + \sqrt 3}, \sqrt 3 \sqrt{1 + \sqrt 3}}
        $$
    \end{enumerate}
\end{ex}

\begin{ex}[H2.5]
    Halla grado y base de $\F_7(t) / \F_7(t^2)$. Halla la expresión de $t^{-1}$ y $(t+1)^{-1}$ en la base que has hallado.\\\\

    Consideramos en polinomio $x^2 - t^2 \in F(t^2)[x]$, donde $t$ es una raíz y $\pm t \notin \F_7(t^2)$. Se puede demostrar por reducción al absurdo que el polinomio $x^2 - t^2$ es irreducible. Por tanto, por el teorema \ref{thm:elemalg}, $|\F_7(t) / \F_7(t^2)| = 2$.\\
    Una $\F_7(t^2)$-base de $\F_7(t)$ es: $\setdef{1, t}$.\\

    Vamos a expresar ahora los elementos que se nos piden. Consideramos $t = f(t)$, $f(x) = x$.
    $$
        x^2 - t^2 = 0,\ x \cdot x = t^2 \implies x \cdot \frac{1}{t^2} \cdot x = 1\text{, que evaluando en $t$: } t \cdot \left(\frac{1}{t^2} \cdot t\right) = 1
    $$
    Con ello, hemos hallado el inverso de $t$. Para hallar el inverso de $t+1$ procedemos de forma parecida. Consideramos $f(x) = x + 1$.
    Vemos que $mcd(f(x), x^2 - t^2) = 1$. Procediendo con el algoritmo de división de polinomios, podemos expresar:
    $$
        x^2 - t^2 = f(x) (x - 1) + (1 - t^2) \implies (x^2 - t^2) + f(x)(1-x) = 1-t^2 \in \F_7(t^2)
    $$
    Entonces:
    $$
        \frac{1}{1 - t^2} (x^2 - t^2) + f(x) \frac{1-x}{1-t^2} = 1
    $$
    Evaluando en $t$:
    $$
        f(t)\frac{1-t}{1-t^2} = 1 \implies (t+1)^{-1} = \frac{1}{1-t^2} \cdot 1 + \frac{1}{t^2 - 1} t
    $$
\end{ex}

\begin{ex}[H2.10]
    Sea $E/K$ una extensión, $\alpha \in E$ algebraico sobre $K$ y $L$ un subcuerpo intermedio. Prueba que $q(x) = Irr(L, \alpha) \divides Irr(K, \alpha) = p(x)$.\\\\

    $p(x) \in K[x] \subseteq L[x]$, entonces $p(x)$ tambien es un polinomio de $L[x]$. Como $p(\alpha) = 0$, por el teorema del elemento algebraico, (teorema \ref{thm:elemalg}) $q(x) \divides p(x)$. Y entonces:
    $$
        |L(\alpha):L| = \delta(q(x)) \leq \delta(p(x)) = |K(\alpha):K|
    $$
\end{ex}

\begin{ex}[H2.11]
    Sea $E/K$ una extensión. Demuestra:
    \begin{enumerate}
        \item Si $|E/K| = p$ con $p$ primo, demuestra que no hay subcuerpos intermedios.\\\\
        Por el teorema de transitividad de grados (teorema \ref{thm:transgr}), sabemos que $|E:K| = |E:L| \cdot |L:K| = p \implies |E:L| = 1 \lor |L:K| = 1$, y por tanto $E = L \lor L = K$.

        \item Sea $|E:K| = p$ con $p$ primo, entonces $E/K$ es simple.\\\\
        Consideramos $L = K(\alpha)$ con $\alpha \in E, \alpha \not in K$, además $K(\alpha) \neq K$ y por tanto (por el primer apartado) $L = E = K(\alpha)$ y es simple ya que $K(\alpha)$ lo es.

        \item Supongamos $\alpha \in E$ tal que $Irr(K, \alpha) = x^3 + x - 1$. Queremos calcular $Irr(K, \alpha^2)$. Como sabemos que $K(\alpha)/K$ tiene grado $3$, y por el apartado siguiente, $K(\alpha) = K(\alpha^2)$, entonces por el teorema \ref{thm:elemalg}, $\delta(Irr(K, \alpha^2))=2$.\\\\
        Entonces:
        $$
            \alpha^3 + \alpha^2 = 1 \iff \alpha^6 + 2\alpha^4 + \alpha^2 - 1 = 0
        $$
        Sea $\beta = \alpha^2$, entonces se satisface que:
        $$
            \beta^3 + 2\beta^2 + \beta - 1 = 0
        $$
        y con ello hemos hallado el polinomio irreducible que buscábamos.

        \item Si $\alpha \in E$, con $|K(\alpha):K| = n$ impar, calcula $|K(\alpha^2):K|$.\\\\
        Como $n$ es impar sabemos que $\alpha^2 \in K(\alpha)$ y por tanto $K(\alpha^2) = K(\alpha)$.\\

        \item Sea $K \subseteq L_1, L_2 \subseteq E$ dos cuerpos intermedios de grado coprimos sobre $K$ demuestra que $L_1 \cap L_2 = K$.\\\\
        Se considera el cuerpo $L_1 \cap L_2$. Sea $d = |L_1 \cap L_2 / K$, $n = |L_1/K|$ y $m=|L_2/K$. Entonces por el teorema \ref{thm:transgr}, $d \divides n$ y $d \divides m$ con $n, m$ coprimos, por tanto $d = 1$ y $\L_1 \cap L_2 = K$.
    \end{enumerate}
\end{ex}

\begin{ex}[H2.12]
    Sea $E/K$ una extensión, y sean $a, b \in E$ algebraicos con $|K(a):K| = n$, $|K(b):K| = m$. Prueba que:

    \begin{enumerate}
        \item $|K(a, b) : K(b) \leq n$.\\\\
        Mirar el ejercicio H2.10

        \item Sean $n$ y $m$ son coprimos, entonces $K(a)\cup K(b) = K$ y $|K(a, b):K|=nm$. Deduce que $Irr(K, a) = Irr(K(b), a)$.\\\\
        Mirar ejercicio H2.11\\
        \item Calcula $Irr(\Q, \alpha)$ con $\alpha = \sqrt{3} + \sqrt[3]{2}$
    \end{enumerate}
\end{ex}

% \part{Segundo parcial}
% % !TeX root = ../galois.tex

\chapter{Extensiones de Galois}

\section{Cuerpos de escisión}

\begin{dfn}[Escisión de un polinomio]
    Sea $f \in K[x]$, $K$ cuerpo, decimos que $f$ se \textbf{escinde} si $\exists a, a_i\in K$ si:
    $$
        f(x) = a(x - a_1) \ldots (x - a_n)
    $$
\end{dfn}

\begin{obs}
    Si $f(x) \in K[x]$, $\delta f = 1$, entonces $f$ se escinde en $K$.
    $$
        f(x) = ax + b = a \left(x - \left(- \frac{b}{a}\right)\right),\ \ \ a, -\frac{b}{a} \in K
    $$
\end{obs}

\begin{obs}
    Si $f \in K[x]$, se escinde en $E$, sea $E/K$ una extensión:
    $$
        f(x) = a (x - a_i) \ldots (x - a_n) \in E[x]
    $$
    Entonces las raíces de $f$ en cualquier extensión de $E$ son $a_1, \ldots, a_n$
\end{obs}

\begin{dfn}[Cuerpo de escisión]
Sea $f \in K[x]$, entonces $E$ es un cuerpo de escisión de $f$ sobre $K$ si:
    \begin{itemize}
        \item $E/K$ es una extensión.
        \item $f$ se escinde en $E$.
        \item Si $f$ se escinde en $L$, con $L$ un cuerpo intermedio $K \subseteq L \subseteq E$, entonces $L=E$, es decir, $E = K(f)$. (Es el menor subcuerpo con la propiedad de que $f$ se escinda en él).
    \end{itemize}
\end{dfn}

\begin{eg}
    $x^2 + 1$ se escinde en $\C$ pero su cuerpo de escisión es $\Q(i)$.
\end{eg}

\begin{obs}
    Si $f \in K[x]$ se escinde en $E$, entonces $E$ tiene todas las raíces de $f$. Para construir su cuerpo de escisión basta adjuntar todas las raíces al cuerpo sobre el que está definido:
    $$
        f(x) = a (x - a_1) \ldots (x - a_n) \in E[x], \text{ entonces su cuerpo de escisión es } K(f) = K(a_1, \ldots, a_n)
    $$
\end{obs}

\begin{lm}[Escisión de polinomios no constantes]\label{lm:3.1}
    Sea $K$ un cuerpo:\\
    \begin{enumerate}
        \item Si $f \in K[x]$ no constante se escinde en $K$ si y sólo si todos los polinomios en la descomposición en términos de factores irreducibles tienen grado $1$.
        \item Si $f \in K[x]$ no constante que se escinde en $K$ y sea $p \in K[x]$ tal que $p \divides f$, entonces $p$ se escinde en $K$.
    \end{enumerate}
\end{lm}

\begin{proof}
    La demostración se deja como ejercicio.
\end{proof}

\begin{lm}[Cuerpos de escisión y cuerpos intermedios]\label{lm:3.2}
    Sean $E/L$ y $L/K$ extensiones, y $f \in K[x]$ un polinomio no constante. Si $E$ es el cuerpo de escisión de $f$ sobre $K$, entonces $E$ es el cuerpo de escisión de $f$ sobre $L$.
\end{lm}

\begin{proof}
    Como $E = K(f)$, tenemos que $f(x) = a (x - a_1) \ldots (x - a_n) \in E[x]$ y $E = K(a_1, \ldots, a_n)$. Además, $f \in K[x] \subseteq L[x]$. Si el cuerpo de escisión de $f$ sobre $L$ es $E'$, fácilmente llegamos a que $E' = L(f) \subseteq E$ y por tanto: $E = E'$ ya que $E$ era el menor cuerpo de escisión de $f$ sobre $K$.
\end{proof}

Si $f$ es un polinomio sobre $\Q[x]$ de grado $n$, \textit{sabemos} que $\exists \alpha_i \in \C$ con $i \in \setdef{1, \ldots n}$ tal que $f(x) = a (x - \alpha_1) \ldots (x - \alpha_n)$, y hemos visto para calcular su cuerpo de escisión basta con adjuntar las raíces de $f$, es decir, $\Q(f) = \Q(\alpha_1, \ldots, \alpha_n)$.\\

Vamos a ver resultados para generalizar este proceso.

\begin{lm}[Lema de Kronecker]\label{lm:3.3}
    Sea $p \in K[x]$ un polinomio irreducible, entonces existe $E/K$ tal que $p$ tiene una raíz en $E$.
\end{lm}

\begin{proof}
    Sea $L = \quo{K[x]}{\gen{p}}$, ya vimos que $L$ es un cuerpo por ser $p$ irreducible. Además se puede comprobar que:
    $$
        \bar x = x + \gen{p} \in L \text{ es raíz de } \bar p \in \bar K[y]
    $$
    Donde $\bar p = \overline{a_0} + \ldots + \overline{a_n}x^n$ y $K \isom \bar K = \setdef{\bar K \mid k \in K}$.

    Hay que comprobar que $\bar K \subseteq L$. Tras ello, habría que ver que $L = \bar K(\bar x)$.
\end{proof}

\begin{thm}[Existencia de cuerpos de escisión]\label{thm:3.4}
    Sea $K$ un cuerpo, $f \in K[x]$ un polinomio no constante. Entonces existe un cuerpo de escisión $E$ de $f$ sobre $K$.
\end{thm}

\begin{proof}
    Basta construir una extensión en la que $f$ se escinda. La demostración sigue por inducción sobre $\delta(f)$.
    \begin{itemize}
        \item[$\delta(f) = 1$] Entonces por el lema \ref{lm:3.1}, $K$ es un cuerpo de escisión de $f$ sobre $K$.
        \item[$\delta(f)>1$] Sea $p \divides f$ un factor irreducible de $f$ en $K[x]$, por el lema \ref{lm:3.3}, sabemos que existe $E/K$ en la que $p$ tiene una raíz $\alpha \in E$, en particular $f(\alpha) = 0$ y por Ruffini:
        $$
            f(x) = (x - \alpha) g(x),\text{ con } g(x) \in K(\alpha)[x]
        $$
        Como $1 \leq \delta(g) \leq \delta(f) - 1 \leq \delta(f)$, por inducción sabemos que existe un cuerpo de escisión $L$ de $g$ sobre $K(\alpha)$. Si $\alpha_1, \ldots \alpha_n \in L$ son las raíces de $g$, entonces:
        $$
            L = K(\alpha)(\alpha_1, \ldots \alpha_n) = K(\alpha, \alpha_1, \ldots, \alpha_n)
        $$
        y por tanto $f(x) = (x - \alpha) g(x) = (x - \alpha) p (x - \alpha_1) \ldots (x - \alpha_n) \in L[x]$, así que $f$ se escinde sobre $L$, y de hecho $L$ es un cuerpo de escisión de $f$ sobre $K$.
    \end{itemize}
\end{proof}

\begin{ex}[Cálculo del cuerpo de escisión]
    Describir el cuerpo de escisión de $f(x) = x^4 - 4x^2 + 2$ sobre $\Q$.\\\\
    Comenzamos hallando las raíces, que resulta fácil pues $f(x) = 0$ es una ecuación bicuadrática. Tenemos entonces las raíces:
    $$
        \alpha = \sqrt{2 + \sqrt 2},\ \beta = \sqrt{2 - \sqrt{2}},\ -\alpha,\ -\beta
    $$
    De donde vemos que $f$ se escinde en $R$ ya que todas las raíces son reales. Para ver su cuerpo de escisión vamos a adjuntar las raíces a $\Q$.
    $$
        E = \Q(\alpha,  \beta, -\alpha, -\beta) = \Q(\alpha, \beta)
    $$
    Sin embargo, podemos reducirlo aún más si vemos que $\Q(\alpha, \beta) = \Q(\beta)$.
    $$
        \alpha\beta = \sqrt{(2 + \sqrt 2)(2 - \sqrt 2)} = \sqrt 2,\ \alpha^{-1} = \left(\frac{\beta}{\sqrt 2}\right) \in \Q(\alpha) \implies \beta = \frac{\beta}{\sqrt 2} \cdot \sqrt 2 \in \Q(\alpha)
    $$
    Con lo que el cuerpo de escisión es: $\Q(\sqrt{2 + \sqrt 2})$.
\end{ex}

\begin{obs}
    Si $E = K(f)$, entonces $E/K$ es finita por el teorema \ref{thm:extalg}. Anteriormente, vimos que si $f \in K[x]$, entonces existe un cuerpo de escisión de $f$ sobre $K$.\\
    Sea $p$ un polinomio irreducible tal que $p \divides f$ y $\delta p > 1$, entonces consideramos $\quo{K[x]}{\gen{p}}$ y en ese cuerpo hay una raíz de $p$ (que también es raíz de $f$), y después aplicamos inducción.
\end{obs}

\begin{thm}\label{thm:3.5}
    Sea $K_1, K_2$ cuerpos, y $\sigma: K_1 \to K_2$ un isomorfismo de cuerpos. Sea $f_1 \in K_1[x]$ no constante y $\delta(f_1) = f_2$.\\ Supongamos que $E_i$ son cuerpos de escisión de $f_1$ sobre $K_i$ con $i \in \setdef{1, 2}$. Entonces existe un isomorfismo $\tau: E_1 \to E_2$ que extiende a $\sigma$, es decir, $\left. \tau \right|_{K_1} = \sigma$.
\end{thm}

\begin{proof}
    Pedir a Santorum.
\end{proof}

\begin{cor}[Unicidad de los cuerpos de escisión]
    Sea $K$ un cuerpo y $f \in K[x]$ no constante. Si $E_1$ y $E_2$ son cuerpos de escisión de $f$ sobre $K$ entonces $\exists \tau$ un isomorfismo $\tau: E_1 \to E_2$ tal que $\tau(a) = a$ $\forall a \in K$.
\end{cor}
\begin{proof}
    Basta tomar $\sigma = id_K$ en el teorema \ref{thm:3.5}.
\end{proof}

%% TODO: COMPLETAR EL DIA 22/10

%%%%%%%%%%%%%%%%%%%%%%%%%%%%%%%%%%%%%%%%%%%%%%%%%%%%%%%%%%%%%%%%%%%%%%%%%% 24/10

\begin{lm}[Normalidad de extensiones en cuerpos intermedios]\label{lm:3.7}
    $K \subseteq L \subseteq E$ extensiones, entonces:
    $$
        E/K \text{ normal } \implies E/L \text{ normal.}
    $$
\end{lm}

\begin{proof}
    Es consecuencia directa del lema \ref{lm:3.2} y de la definición.
\end{proof}


\begin{cor}[al teorema \ref{thm:3.5}]
    Sea $E/K$ una extensión normal, $K \subseteq M_1, M_2 \subseteq E$ con $M_1, M_2$ cuerpos intermedios.\\\\
    Si existe un isomorfismo $\sigma: M_1 \to M_2$ que fija a $K$, es decir, $\sigma(a) = a,\ \forall a \in K$, entonces $\sigma$ se extiende a un isomorfismo $\tau: E \to E$ ($\left. \tau \right|_{M_1} = \sigma$).
\end{cor}
\begin{proof}
    Sea $E = K(f)$ con $f \in K[x]$ no constante. En particular por el lema \ref{lm:3.2}, $E = M_1(f) = M_2(f)$. Como $\sigma(f) = f$, el resultado se sigue del teorema \ref{thm:3.5}.
\end{proof}
\begin{obs}
    La hipótesis $\sigma(a) = a,\ \forall a \in K$ es necesaria. La idea es que sin esta hipótesis:
    $$
        E = M_1(f) = M_2(f)
    $$
    pero no sabríamos si $E = M_2(\sigma(f))$.
\end{obs}

\begin{thm}[Condición necesaria y suficiente para la normalidad de una extensión]\label{thm:3.9}
    Sea $E/K$ una extensión finita, $E/K$ es normal si y sólo si todo polinomio $p \in K[x]$ irreducible con una raíz en $E$ se escinde en $E$.
\end{thm}

\begin{proof}$ $
    \begin{itemize}
        \item[$\ \Longleftarrow\ $] Sabemos que $E/K$ es finita, digamos $|E:K| = n$, podemos tomar una $K$-base $\sdf{a_1, \ldots, a_n}$. En particular $E=K(a_1, \ldots, a_n)$.\\
        Como $E/K$ es finita, por el teorema \ref{thm:2.2}, los $a_i$ en son algebraicos sobre $K$. Sea $p_i = Irr(k, a_i) \in K[x]$, por hipótesis cada $p_i$ se escinde en $E$.\\

        Tomamos $f = \prod_{i=1}^{n} p_i \in K[x]$ se escinde en $E$. De hecho, $E = K(f)$. De esta forma hemos probado que:
        $$
            K(f) \subseteq E = K(a_1, \ldots, a_n) \subseteq K(f)
        $$
        \item[$\implies$] Ahora partimos de que $E/K$ es finita y normal. Entonces $\exists f \in K[x]$ tal que $E = K(f)$.\\
        De hecho, si $f(x) = c(x - b_1) \ldots (x-b_n) \in E[x]$, $E = K(b_1, \ldots, b_n)$.\\
        Sea $p\in K[x]$ irreducible con una raíz $a\in E$, por Ruffini (corolario \ref{cor:ruffini}) $p(x) = (x-a)g(x)\in E[x]$. Además, por el teorema \ref{thm:3.4}, existe $M/E$ un cuerpo de escisión de $p$ sobre $E$, por tanto:
        $$
            p(x) = d(x-d_1) \cdot \ldots \cdot (x-d_r), \text{ podemos suponer que $a = d_1$}
        $$
        Dado $d_i$ con $i > 1$, queremos ver que $d_i \in E$. Llamemos $b = d_i$ para algún $i$. Como $a$ y $b$ son raíces de $p \in K[x]$ irreducible, por el teorema \ref{thm:2.5} $\exists \sigma: K(a) \to K(b)$ tal que $\sigma(k) = k,\ \forall k \in K$.

        %
        %\sigma: K(a) --------> K(b)
        %         ^              ^
        %         |              |
        %         v              v
        %id:      K  --------->  K

        Como $K(a) \subseteq E$, en particular $E = K(a)(f)$. Por otro lado $K(b)(f) = E(b)$. Las raíces de $f$ en $E(b)$ son $b_1, \ldots, b_m$, así que:
        $$
            K(b)(f) = K(b)(b_1, \ldots, b_m) = K(b_1, \ldots, b_m) = E(b)
        $$
        Por el teorema \ref{thm:3.5} existe un isomorfismo $\tau: E \to E(b)$ que extiende a $\sigma$.
        %
        %\tau:  E=K(a)(f) ----> E(b) = K(b)(f)
        %         |              |
        %         |              |
        %\sigma: K(a) --------> K(b)
        %         ^              ^
        %         |              |
        %         v              v
        %id:      K  --------->  K

        Por el ejercicio $H2.15$, $|E:K| = |E(b):K|$ y $|K(a):K| = |K(b):K|$. Además $|E:K(a)| = |E(b):K(b)|$.
        Por tanto:
        $$
            |E:K| = |E:K(a)||K(a):K| = |E(b):K(b)||K(b):K| = |E(b):K| = |E(b):E||E:K| \text{ usando el teorema \ref{thm:transgr}.}
        $$
        Con lo que concluimos:
        $$
            |E(b):E| = 1 \implies E(b) = E \implies b\in E
        $$
    \end{itemize}
\end{proof}


\begin{ex}[H3.5]
    Decide si son normales las siguientes extensiones:
    \begin{enumerate}
        \item $\Q(\sqrt{5}i)/\Q$.\\\\

        Es normal ya que $\Q(\sqrt{5}i) = \Q(x^2+5)$ es el cuerpo de escisión de $x^2+5$.
        \item $\Q(\sqrt{5})/\Q$.\\\\

        Es normal ya que $\Q(\sqrt{5})/\Q = \Q(x^2-5)$ es el cuerpo de escisión de $x^2-5$.

        \item $\Q(\sqrt[4]{5}/\Q)$.\\\\

        No es normal. Sabemos que $\Q(\sqrt[4]{5})/\Q$ es finita, y como $p(x) = x^4 - 5$ es irreducible en $\Q[x]$, tiene una raíz en $\Q(\sqrt[4]{5})$ pero no escinde, por la caracterización de normalidad (teorema \ref{thm:3.9}), concluimos con que $\Q(\sqrt[4]{5})/\Q$ no es normal.
    \end{enumerate}
\end{ex}

\begin{ex}[H3.6]
    Demuestra que $\Q(\sqrt[3]{2})/\Q$ no es normal. Encuentra una extensión normal que contenga $\Q(\sqrt[3]{2})$ como subcuerpo.\\\\

    Sabemos que no es normal porque sea $p(x) = x^3 - 2 \in \Q[x]$, se puede comprobar que $p(x)$ es irreducible y tiene una raíz en $\Q(\sqrt[3]{2})$ pero no se escinde. Por tanto, por el teorema \ref{thm:3.9}, el cuerpo no es normal.\\
    Una extensión normal que contiene a $\Q(\sqrt[3]{2})$ puede ser $\Q(\sqrt[3]{2}, \omega)$ con $\omega \in \C,\ o(\omega) = 3$.
\end{ex}

\begin{obs}
    Si $E/K$ es normal y $f\in K[x]$ no necesariamente irreducible que tiene una raíz en $E$, $f$ no tiene por que escindirse.
\end{obs}

% \part{Tercer parcial}
% % !TeX root = ../galois.tex

\chapter{Teoría de Galois}

\section{Teorema Fundamental} % TODO: Comprobar si es la primera seccion

A partir de este capítulo vamos a trabajar siempre con extensiones $E/K$ que sean de Galois, es decir, $E/K$ es separable y normal. En el teorema \ref{thm:3.11}, caracterizamos al grupo de Galois de una extensión, vamos a ver que ambas hipótesis (separabilidad y normalidad) son necesarias.

\begin{eg}[Grupo de Galois de una extensión no normal]
    Consideremos $\Q(\sqrt{1 + \sqrt 7})/\Q$. Podemos hallar que $\sqrt{1 + \sqrt 7}$ es raíz de $p(x) = x^4-2x^2-6 \in \Q[x]$, que es irreducible por el criterio de Einsestein para $p=2$.\\
    Entonces, las raíces de $p(x)$ son $\pm\sqrt{1 + \pm\sqrt 7}$, y su cuerpo de escisión es $E = \Q(p) = \Q(\sqrt{1 + \sqrt 7}, \sqrt{1 - \sqrt 7})$.\\
    Si suponemos que $\Q(\sqrt{1 + \sqrt 7}) = E(p)$ rápidamente llegamos a una contradicción con lo que concluimos que $\Q(\sqrt{1 + \sqrt 7})/\Q$ no es normal. Nos  preguntamos ahora acerca de $\gal(\Q(\sqrt{1 + \sqrt 7})/\Q)$.\\
    Por el teorema \ref{thm:3.11}, $\Omega = \sdf{\pm \sqrt{1 + \sqrt 7}}$ es el conjunto de raíces de $p$ en $\Q(\sqrt{1 + \sqrt 7})/\Q$, y por tanto $\gal(\Q(\sqrt{1 + \sqrt 7})/\Q) \leq S_2 \implies \gal(\Q(\sqrt{1 + \sqrt 7})/\Q) = \sdf{id, \sigma}$ con $\sigma(\sqrt{1 + \sqrt 7}) = -\sqrt{1 + \sqrt 7}$.\\\\
    De esta forma vemos que:
    $$
        \vabs{\gal(\Q(\sqrt{1 + \sqrt 7})/\Q)} = 2 \text{ pero } \vabs{\Q(\sqrt{1 + \sqrt 7}):\Q} = 4
    $$
\end{eg}

\begin{ex}[H4.3]$ $
    \begin{itemize}
        \item[(a)] Determina que $\F_3(t) / \F_3(t^3)$ no es de Galois.
        \item[(b)] Determina que $\C(t) / \C(t^3)$ es de Galois.
        \item[(c)] Calcula los grupos de Galois asociados.
    \end{itemize}
    \textbf{Solución}\\
    \begin{itemize}
        \item[(a)] Vamos a ver que la extensión no es separable. Sabemos que $\F_3(t) / \F_3(t^3)$ es algebraica y finita. Consideramos el polinomio irreducible $x^3-t^3 \in \F_3(t^3)[x]$, que tiene $t$ como raíz, entonces no es separable por que $x^3-t^3 = (x-t)^3$ y por tanto no es de Galois.\\
        Además, $\gal(\F_3(t)/\F_3(t^3)) = \sdf{id}$.
        \item[(b)] Sabemos que $\C(t) / \C(t^3)$ es separable por que estamos en un cuerpo de característica cero. Además, sea el polinomio $x^3-t^3 \in \C(t^3)[x]$, sus raíces son $\Omega = \sdf{t, t\omega, t\bar\omega}$ donde $\omega \in \C:\ o(\omega)=3$.\\
        Como $\Omega \subseteq \C(t)$, $\C(t)$ es el cuerpo de escisión de $p(x)$ y por tanto $\C(t)/\C(t^3)$ es separable y normal y por tanto de Galois.
        $$
            \text{Luego, } \vabs{\gal(\C(t) / \C(t^3))} = \vabs{\C(t) : \C(t^3)} = 3
        $$
        Con lo que hallamos:
        $$
            \gal(\C(t) / \C(t^3)) = \gen{\tau}:\ \tau(t) = t\omega,\ o(\tau) = 3
        $$
        \item[(c)] Resuelto en los apartados respectivos.
        \end{itemize}
\end{ex}

\begin{ex}[H4.2]
    Sea $E = \F_{p^n}$, demuestra que $E / \F_p$ es de Galois y calcula $\gal(E / \F_p)$.\\\\

    \textbf{Solución}\\
    Sabemos que $E / \F_p$ es separable por que $\F_p$ es perfecto y la extensión es algebraica. Tenemos que ver que $E / \F_p$ es normal.\\
    Además $E$ es el cuerpo de escisión de $x^{p^n} - x$ sobre $\F_p \implies E / \F_p$ es normal y por tanto es de Galois.\\\\
    Como $E / \F_p$ es de Galois $\implies \vabs{\gal(E / \F_p)} = \vabs{E : \F_p} = n$. Sea $\varphi = Frob \in \gal(E / \F_p)$, sabemos que $o(\varphi) \divides n$.
    Sea $d = o(\varphi)$, sabemos que $d \leq n$ y que:
    $$
        a = id(a) = \varphi^d(a) = a^{p^d} \forall\ a \in E \implies \forall a \in E,\ a^{p^d} - a = 0
    $$
    Por tanto, todo elemento en $E$ es raíz de $x^{p^d} - x \implies E$ es el cuerpo de escisión de $x^{p^d} - x$. Esto quiere decir que $E$ tiene $\delta(x^{p^d} - x)$ elementos, pero $E = \F_{p^n}$ que tiene $p^n$ elementos $\implies d = n$. Por tanto:
    $$
        \gal(E / \F_p) = \gen{\varphi}
    $$
\end{ex}

\begin{cor} % TODO: Dar nombre
    Sea $E/K$ una extensión de Galois y $L$ un cuerpo intermedio ($K \subseteq L \subseteq E$), entonces:
    $$
        L/K \text{ es de Galois } \iff \forall \tau \in \gal(E/K),\ \rest{\tau}_L \in \gal(E/K)
    $$
    Y además, en cualquiera de los supuestos equivalentes anteriores, todo $\sigma \in \gal(L/K)$ se extiende a $\tau \in \gal(E/K)$ (tiene exactamente $\vabs{E:L}$ extensiones), $\gal(E/L) \normsub \gal(E/K)$ y $ \quo{\gal(E/K)}{\gal(E/L)} \isom \gal(L/K)$.\\\\
    Podemos hacer un diagrama:
    % TODO
    % E          {1} = \gal(K/K)
    % |          |
    % |          |
    % L -------- N = \gal(E/L)
    % |          ||
    % |          ||
    % K          G = \gal(L/K)
\end{cor}

El objetivo de este capítulo es probar el teorema \ref{thm:fund-gal}:

Para ello vamos a hacer uso del corolario \ref{cor:3.12}. Además, necesitaremos un resultado sobre automorfismos de cuerpos en general.\\
Recordemos que si $E$ es un cuerpo, entonces $\aut(E)$ actúa sobre $E$ según:
$$
    \forall e \in E,\ \forall \sigma \in \aut(E),\ \ e \cdot \sigma = \sigma(e)
$$
y cumple las propiedades:
\begin{align*}
    e \cdot (\sigma \tau) = (e \cdot \sigma) \tau\\
    e \cdot (\sigma \tau) = \sigma \circ \tau (e) = \tau ( \sigma (e)) = \tau(e \cdot \sigma) = (e \cdot \sigma) \tau
\end{align*}
Además, si $G \leq \aut(E)$, entonces $G$ actúa sobre $E$, y recordemos la notación del conjunto de puntos fijos de $E$ por la acción de $G$:
$$
    E^G = \sdf{e \in E \mid e \cdot \sigma = e} = \sdf{e \in E \mid \sigma(e) = e} \subseteq E \text{ es un subcuerpo }
$$
Por tanto $E/E^G$ es una extensión.

\begin{thm}[de Artin]\label{thm:4.1}
    Sea $E$ un cuerpo, y sea $G \leq \aut(E)$ un subgrupo finito, entonces:
    $$
        \vabs{E:E^G} = \vabs{G}
    $$
\end{thm}
\begin{proof}
    Se proporciona un esquema de la prueba.\\
    Sea $F = E^G$, la prueba sigue por demostrar la doble igualdad:
    \begin{enumerate}
        \item $m = \vabs{E:F} \geq \vabs{G} = n$.\\
        Por reducción al absurdo suponemos $m < n$, entonces podemos encontrar $\tau_1, \ldots, \tau_n \in G$ distintos tal que:
        $$
            \sum_{j=1}^n a_j\tau_j = 0
        $$
        Y esto entra en contradicción con el teorema de Dedekind (teorema \ref{thm:4.2}). Entonces $m \geq n$.
        \item $\vabs{E:F} = m \leq n = \vabs{G}$\\
        Tendríamos que trabajar con sistemas lineales, soluciones y manipulaciones...
    \end{enumerate}
\end{proof}

\begin{thm}[de Dedekind]\label{thm:4.2}
    Sea $E$ un cuerpo y sean $\tau_1, \ldots, \tau_n \in \aut(E)$ distintos. Entonces $\tau_1, \ldots, \tau_n$ son $E-$linealmente independientes (en el espacio vectoreial de funciones de $E$).
\end{thm}
\begin{proof}
    Queremos probar que si $a_j \in E$:
    $$
        \sum_{j=1}^{n} a_j\tau_j = 0 \implies a_j = 0 \forall j
    $$
    Procedemos por inducción sonre $n$.
    \begin{itemize}
        \item[$n=1$]
        $$
            a_1\tau_1 = 0 \implies a_1\tau_1(1) = a_1 = 0
        $$
        \item[$n>1$]
        \begin{align}
            \forall x  \in E:\ \sum_{j=1}^n a_j\tau_j(x) = 0\\
            \forall y \in E:\ \sum_{j=1}^n a_j\tau_j(xy) = \sum_{j=1}^n a_j\tau_j(x)\tau_j(y) = 0\\
            \forall y \in E:\ \sum_{j=1}^n a_j\tau_j(x)\tau_n(y) = 0
        \end{align}
        Entonces restando:
        $$
            0 = \sum_{j=1}^n a_j(\tau_j(y) - \tau_n(y)) \tau_j(x) = \sum_{j=1}^{n-1} a_j(\tau_j(y)-\tau_n(y))\tau_j(x)
        $$
        Por inducción:
        $$
            a_j(\tau_j(y) - \tau_n(y)) = 0 \forall j \in \sdf{1, \ldots, n-1}
        $$
        Supongamos que $\exists a_j \neq 0$, entonces:
        $$
            \tau_j(y) = \tau_n(y) \implies \tau_j = \tau_n \text{ y es imposible por hipótesis.}
        $$
    \end{itemize}
\end{proof}

\begin{thm}[de puntos fijados]\label{thm:4.3}
    Sea $E/K$ una extensión de Galois, y $G = \gal(E/K)$ su grupo de Galois $\implies E^G = K$.\\
    Es decir, si $a \in E$ tal que $\forall \tau \in G:\ \tau(a) = a \implies a \in K$.
\end{thm}
\begin{proof}
    Escribimos $F = E^G \subseteq E$. Suponemos $K \subseteq F$, tenemos entonces el diagrama:
    %TODO
    %   |-- E ----------|
    %   |   |           | |G|
    %|G||   F = E^G ----|
    %   |   |           | 1
    %   |-- K ----------|
    Por el teorema \ref{thm:3.4.6}, sabemos que $\vabs{G} = \vabs{E:K}$. Por el teorema de Artin (teorema \ref{thm:4.1}), $\vabs{E:E^G} = \vabs{E:F} \implies |F:K| = 1 \implies K = F = E^G$.
\end{proof}
\begin{obs}
    El recíproco también es cierto, se corresponde con el ejercicio H4.4. Basta ver que $\forall \alpha \in E,\ Irr(K, \alpha)$ se escinde en $E$ y todas sus raíces son distintas.
\end{obs}
\begin{cor}\label{cor:4.4}%TODO Dar nombre
    Sea $E/K$ una extensión de Galois, $G = \gal(E/K)$ su grupo de Galois y $L$ un cuerpo intermedio entonces:
    $$
        L/K \text{ es normal } \iff \gal(E/L) \normsub G
    $$
\end{cor}
\begin{proof}$ $
    \begin{itemize}
        \item[($\implies$)] Por el corolario \ref{cor:3.12}(b), entonces:
        $$
            (\star)\ L/K \text{ normal } \iff \forall \tau \in G: \tau(L) = L
        $$

        \item[($\Longleftarrow$)] Partimos de $\gal(E/L) \normsub G$, y queremos ver que $L/K$ es normal. Usando $(\star)$, basta ver que $\forall \tau \in G$, $\tau(L) = L$. Entonces:
        $$
            H \normsub G \iff \forall \sigma \in H,\ \forall \tau \in G,\ \sigma^\tau = \tau^{-1} \sigma \tau \in H
        $$
        entonces $\forall l \in L$:
        \begin{align*}
            \sigma^\tau(l) = \tau^{-1} \sigma \tau(l) = l\cdot \tau^{-1} \cdot \sigma \cdot \tau = l\\
            l \cdot (\tau^{-1} \cdot \sigma) = l \cdot \tau^{-1} \implies \tau^{-1}(l) \in E^H = L
        \end{align*}
        Y por tanto $E/L$ es de Galois por el teorema \ref{thm:4.3}. Como $E/L$ es de Galois $\implies \tau^{-1}(L) \subseteq L \implies L \subseteq \tau(L)$.\\
        Entonces, tenemos el diagrama:
        %   tau rest L : L \to (isom) tau(L)
        %           |                   | 1
        %           |                   |L cuerpo intermedio
        %           K \to ------------> K
        %
        Y por el ejercicio H2.15, $\vabs{L:K} = \vabs{\tau(L):K} = \vabs{\tau(L):L}\vabs{L:K} \implies \tau(L) = L$.
        Además, también tenemos el diagrama:
        %        E
        %      /   \
        %  \tau(L)  L
        %      \   /
        %        K
        Entonces:
        \begin{align*}
            H = \gal(E/L) \leq G = \gal(E/K),\ \tau \in G\\
            H \isom H^\tau = \tau^{-1} H \tau = \sdf{\tau^{-1} \sigma \tau \mid \sigma \in H}\\
            H^\tau = \gal(E/\tau(L))
        \end{align*}
    \end{itemize}
\end{proof}

\begin{obs}[Notación]
    Recordamos que usamos la notación de acción de grupos a derechas, es decir:\\
    Sean $f,\ g$ funciones $f,g: E \to E$ con $E$ un cuerpo, $f \circ g (e) = g (f (e))$ y sea el grupo de automorfismos $(\aut(E), \circ)$, entonces para $\sigma, \tau \in  \aut(E)$ tenemos que $\forall e \in E$:
    $$
        e \cdot \sigma = \sigma(e),\text{ y } e \cdot (\sigma \tau) = (e \cdot \sigma) \tau
    $$
\end{obs}

\begin{lm}[Lema]\label{lm:4.5} %%TODO: DAR NOMBRE
    Sean $K \subseteq L \subseteq E$ cuerpos, $H = \gal(E/L) \leq \gal(E/K)$, y sea $\tau \in G$ entonces:
    $$
        H^\tau = \gal(E/\tau(L))
    $$
\end{lm}
\begin{proof}$ $
    \begin{align*}
        %% TODO: Pedir a Santorum
    \end{align*}
\end{proof}


\begin{thm}[Fundamental de la Teoría de Galois]\label{thm:fund-gal}\label{thm:4.6}
    Sea $E/K$ una extensión de Galois con $G = \gal(E/K)$, y sean:
    $$
        \mathcal{K} = \sdf{L \mid K \subseteq L \subseteq E},\ \mathcal{J} = \sdf{H \mid H \leq G}
    $$
    Entonces:
    \begin{itemize}
        \item[(a)] Las aplicaciones:
        $$ f: \mathcal{K} \to \mtc{J};\ L \mapsto \gal(E/L)$$
        $$ g: \mtc{J} \to \mtc{K};\ H \mapsto E^H $$
        son biyectivas y $f^{-1} = g,\ g^{-1} = f$.
        \item[(b)] Bajo esta correspondencia:
        $$
            L/K \text{ es normal } \iff \gal(E/L) \normsub G
        $$
        \item[(c)] Si $L \longleftrightarrow H$, $\vabs{L:K}$ (grado) $ = \vabs{G:H}$ (índice)
    \end{itemize}
\end{thm}
\begin{proof}$ $
    \begin{itemize}
        \item[(a)] Sea $L$ un cuerpo intermedio, como $E/L$ es una extensión de Galois, si $H = f(L) = \gal(E/L)$ entonces:
        $$
            (f \circ g)(L) = g(f(L)) = g(H) = E^H =^{\text{(teorema \ref{thm:4.3})}} = L \implies (f \circ g) = id_{\mtc K}
        $$
        En particular, $g$ es sobreyectiva, si probamos que es inyectiva, entonces $g^{-1} = f$.\\
        Sean $H_1, H_2 \normsub G$ tales que:
        $$
            L = E^{H_1} = g(H_1) = g(H_2) = E^{H_2}
        $$
        Por el teorema de Artin (teorema \ref{thm:4.1}), entonces:
        \begin{align*}
            \vabs{E:E^{H_1}} &= \vabs{H_1}\\
            \vabs{E:E^{H_1}} &= \vabs{E:L} = \vabs{E:E^{H_2}} = \vabs{H_2}\\
            \vabs{E:L} &= \vabs{E:E^{\gal(E/L)}} = \vabs{\gal(E/L)}
        \end{align*}
        Como $E/L$ es de Galois, entonces
        % TODO: Pedir a Santorum
    \end{itemize}
\end{proof}

\begin{ex}[H4.5]
    Calcular el grupo de Galois del cuerpo de escisión de $x^12 - 1$ sobre $\Q$ y todas las subextensiones intermedias.\\\\
    \textbf{Solución}\\
    Comenzamos viendo que:
    \begin{align*}
        E &= \Q(x^12-1) = \Q(x^6+1)\\
          &=\Q(\pm 1, \pm i \pm \frac{\sqrt 3}{2} \pm \frac{1}{2}i, \pm\frac{1}{2} + \frac{\sqrt 3}{2})\\
          &= \Q(\sqrt 3, i)
    \end{align*}
    Además, podemos comprobar que $\vabs{E:\Q} = 4 = \vabs{G}$ y como $\vabs{G} = 2^m$, entonces $G = \gal(E/\Q)$ es normal.\\
    Recordamos que si un grupo $G$ es cíclico, entonces $\forall d \divides \vabs{G}$ existe un único subgrupo de índice $d$.\\
    Además, $G$ tiene dos grupos distintos de índice $2 \implies G = C_2 \times C_2$.
    Por el teorema \ref{thm:4.6}, existen tantas subextensiones como $H \leq G \isom C_2 \times C_2$. Es decir, tantos como subgrupos tiene $C_2 \times C_2$, que podemos calcularlos y hallar que son $5$.
\end{ex}
\begin{obs}
    Si nos fijamos en la prueba del teorema \ref{thm:3.5}, entonces:
    %
    % K_1(f_1) = E_1 ---------> E_2 = K_2(f_2)
    %          |                          |
    %          |                          |
    %          |                          |
    %  \sigma: K_1 ---------------------> K_2
    %
    Sea $f_2 = \sigma^\star(f_1)$, entonces $\sigma$ se puede extender a $\tau: E_1 \to E_2$. Para cada $p_1 \divides f_1$ factores irreducibles de $f_1$ y para cada par $\alpha$ raíz de $p_1$ en $E_1$ y $\beta$ raíz $\sigma^\star(p_1) = p_2 \divides f_2$ existe $\tau: E_1 \to E_2$ tal que:
    $$
        \tau(\alpha) = \beta
    $$.
\end{obs}

%%%%%%%%%%%%%%%%%%%%%%%%%%%%%%%%%%%%%%%%%%%%%%%%%%%%%%%%%%%%%%%%%%%%%%%%%% 26/11
%TODO: Completar el ejercicio H4.5

Vamos a recordar algunos resultados de teoría de grupos.
\begin{pro}[Resultados de teoría de grupos]$ $
    \begin{enumerate}
        \item Si $G$ es cíclico $\implies G$ es abeliano.
        \item Si $\vabs{G} = p^a$ con $p$ un primo, entonces $Z(G) > 1$.
        \item Si $\vabs{G} = p^2$ con $p$ un primo, entonces $G$ es abeliano, en particular $G \isom C_{p^2}$ o $G \isom C_p \times C_p$.
    \end{enumerate}
\end{pro}

\begin{eg}[Cálculo del grupo de Galois de una extensión dado un polinomio (1)]
    Sea $x^4-2$ un polinomio en $\Q[x]$, calcula $\gal(E/\Q)$ siendo $E$ su cuerpo de escisión.\\\\

    \textbf{Solución}\\
    Las raíces en $\C$ del polinomio son $\sdf{\pm \sqrt[4]{2}, \pm \sqrt[4]{2} i}$, por tanto $E = \Q(\pm \sqrt[4]{2}, \pm \sqrt[4]{2} i)  = \Q(\sqrt[4]{2}, i)$.
    Ahora calculamos el grado de la extensión, por transitividad de grados: $\vabs{E:\Q(\sqrt[4]{2})} = 2$ y $\vabs{\Q(\sqrt[4]{2}):\Q} = 4$ por tanto $\vabs{G} = \vabs{E:\Q} = 8$, con $G = \gal(E/\Q)$.\\

    Además, sabemos que $\Q(\sqrt[4]{2})/\Q$ no es normal, lo que implica que $G$ tiene un subgrupo que no es normal (por el teorema \ref{thm:4.6}) y entonces $G$ no puede ser abeliano. De hecho $\gal(E/\Q(\sqrt[4]{2})) = \sdf{\1_E, \tau \mid \tau(i) = -i}$. Por tanto $G \isom D_8$ o $G \isom Q_8$.
    Sin embargo, en $Q_8$ todo subgrupo es normal, por lo que $G \isom D_8$.\\

    Recordamos que $D_8 = \gen{a, b\ \mid a^4 = b^2 = 1,\ a^b = b^{-1} a b = a^{-1}}$. En general:
    $$
        D_{2n} = \gen{\alpha, \beta\ \mid a^n = 1 = \beta^2,\ \alpha^\beta = \alpha^{-1}}
    $$
    Para calcular el grupo de Galois vamos a considerar $L = \Q(\sqrt 2, i)$, y ya hemos visto que $L = \gal(L/\Q) \isom C_2 \times C_2 = \sdf{\1_L, \alpha_2, \alpha_3, \alpha_4}$.\\

    \begin{center}
        \begin{tabular}{|c||c|c|c|}
        \hline
                   & $\sqrt{2}$  & $i$  & ord \\ \hline \hline
        $id$       & $\sqrt{2}$  & $i$  & $1$ \\ \hline
        $\alpha_2$ & $\sqrt{2}$  & $-i$ & $2$ \\ \hline
        $\alpha_3$ & $-\sqrt{2}$ & $i$  & $2$ \\ \hline
        $\alpha_4$ & $-\sqrt{2}$ & $-i$ & $2$ \\ \hline
        \end{tabular}
    \end{center}

    Entonces sea $E = L(x^2-\sqrt{2})$, sabemos que $\alpha_i: L \to L$ es un ismorfismo, y que su extensión $\hat{\alpha}: E \to E$, puede llevar raíces de $x^2 - \sqrt{2}$ (si $\alpha(\sqrt 2) = \sqrt 2)$ o raíces de $x^2 - \sqrt{2}$ en raíces de $x^2 + \sqrt{2}$ (si$ \alpha(\sqrt 2) = -\sqrt 2$).\\

    Además ya hemos comentado que $G \isom D_8$. Podemos hacer una tabla de sus elementos:\\
    \begin{center}
        \begin{tabular}{|c||c|c|c|}
        \hline
                 & $\sqrt[4]{2}$   & $i$  & ord \\ \hline \hline
        $\1_E$     & $\sqrt[4]{2}$   & $i$  & $1$ \\ \hline
        $\tau_2$ & $-\sqrt[4]{2}$  & $i$  & $2$ \\ \hline
        $\tau_3$ & $\sqrt[4]{2}$   & $-i$ & $2$ \\ \hline
        $\tau_4$ & $-\sqrt[4]{2}$  & $-i$ & $2$ \\ \hline
        $\tau_5$ & $\sqrt[4]{2}$   & $i$  & $4$ \\ \hline
        $\tau_6$ & $-\sqrt[4]{2}i$ & $i$  & $4$ \\ \hline
        $\tau_7$ & $\sqrt[4]{2}$   & $-i$ & $2$ \\ \hline
        $\tau_8$ & $-\sqrt[4]{2}i$ & $-i$ & $2$ \\ \hline
        \end{tabular}
    \end{center}

Queremos encontrar una presentación del grupo. Sabemos que $\tau_5^{-1}$ tiene orden $4$ distinto de $\tau_5$, por tanto $\tau_5^{-1} = \tau_6$. Por otra parte:
\begin{align*}
    &o(\tau_3) = 2 \iff \tau_3^2 = \1_E \iff \tau_3^{-1} = \tau_3\\
    &\tau_5^{\tau_3} = \tau_3^{-1}\tau_5\tau_3 = \tau_3\tau_5\tau_3 \implies \\
    &\tau_5^{\tau_3}(i) = i,\ \tau_5^{\tau_3}(\sqrt[4]{2}) = -\sqrt[4]{2}i \implies\\
    &\tau_5^{\tau_3} = \tau_6 = \tau_5^{-1}
\end{align*}
Por tanto:
$$
    G = \gen{\tau_5,\ \tau_3 \mid \tau_5^4 = \tau_3^2 = 1,\ \tau_5^{\tau_3} = \tau_5^{-1} }
$$
\end{eg}
\begin{ex}[H4.7]
    Sea $p$ un primo, y $f(x) = x^p - 1 \in \Q[x]$.
    \begin{itemize}
        \item[(a)] Halla $E = \Q(f)$
        \item[(b)] Prueba que $E/\Q$ es simple de grado $p - 1$
        \item[(c)] Demuestra que $G = \gal(E/\Q)$ es cíclico encontrando un isomorfismo explícito entre $G$ y $F_p^\times$
        \item[(d)] Demuestra que si $p$ es impar, entonces $E$ contiene exactamente una extensión cuadrática de $\Q$ (es decir, una extensión de grado $2$ sobre $\Q$).
    \end{itemize}

    \textbf{Solución}\\\\
    \begin{itemize}
        \item[(a)] Las raíces de $f$ son las raíces $p$-ésimas de la unidad, por tanto, sea $\Omega$ el conjunto de raíces:
        $$
            \Omega = \sdf{1, \xi, \ldots, \xi^{p-1}},\ \xi = e^{\frac{2\pi i}{p}} \implies E = \Q(\xi)
        $$
        \item[(b)] $\Q(\xi)/\Q$ es simple pues ya hemos visto que $\xi$ la genera. Además, como $f$ es ciclotómico:
        $$
            \vabs{E:\Q} = \delta(Irr(\Q, \xi)) = \delta(x^{p-1} + \cdots + x + 1) = p-1
        $$
        \item[(c)] $\sigma \in \gal(\Q(\xi)/\Q)$ queda determinada por $\sigma(\xi)$. Además sabemos que $\vabs{\gal(\Q(\xi)/\Q)} = p-1$. Por tanto:\\
        \begin{center}
            \begin{tabular}{|c|c|}
            \hline
            $G = \gal(\Q(\xi)/\Q) $ & $ \xi $\\ \hline \hline
            $\1 = \sigma_1 $ & $ \xi $\\ \hline
            $\sigma_2 $ & $ \xi^2 $\\ \hline
            $\vdots $ & $ \vdots $\\ \hline
            $\sigma_{p-1} $ & $ \xi^{p-1} $\\ \hline
            \end{tabular}
        \end{center}
        Tenemos que encontrar un isomorfismo entre $\F_p^\times$ y $G$. Recordemos que $F_p^\times = \sdf{1, 2, \ldots, p-1}$. Entonces, sea $\rho: G \to F_p^\times$ nuestro isomorfismo, basta establecer $\rho(\sigma_j) = j$, entonces:
        $$
            \rho(\sigma_k \sigma_j) = \rho(\sigma_{kj}) = \overline{kj} = \rho(\sigma_k) \rho(\sigma_j)
        $$
        Y es fácil ver que $\rho$ es biyectiva. Con esto hemos hallado que $G \isom \F_p^\times \isom C_{p-1}$.

        \item[(d)]  Por el TFTG (\ref{thm:fund-gal}), $G$ tiene un único subgrupo $H$ con $\vabs{G:H} = 2 \iff $ se corresponde con la única subextensión de $\Q(\xi) / \Q$ de grado $2$ sobre $\Q$. Es decir:
        $$
            \Q(\xi)^H = \Q(\xi)^{\sigma_j} \text{ (habría que calcular un elemento $\sigma_j$ tal que $o(\sigma_j) = \frac{p-1}{2}$)}
        $$
    \end{itemize}
\end{ex}



\begin{pro}[Teoremas de Sylow]
    Sea $G$ un grupo de orden $\vabs{G} = p^a n$ con $p$ primo y $p \not \divides n$, entonces:
    \begin{enumerate}
        \item $\exists P \leq G$ subgrupo, de orden $\vabs{P} = p^a$. Decimos que $P \in Syl_p(G)$.
        \item Sean $P_1, P_2 \in Syl_p(G) \implies \exists g \in G: P_1^g = P_2$, es decir, $\exists g \in G: \sdf{g^{-1} \cdot x \cdot g \mid x \in P_1} = P_2$.
        \item $\vabs{Syl_p(G)} = \vabs{G:N_G(P)}$.
    \end{enumerate}
\end{pro}

\begin{eg}[Cálculo del grupo de Galois de una extensión dado un polinomio (2)]
    Partimos del polinomio: $f = x^3 - 2 \in \Q[x]$, vamos a hallar el grupo de Galois de su cuerpo de escisión.
    \begin{enumerate}
        \item Calculamos $E = \Q(f)$.
        $$
            E = \Q(\sqrt[3]{2}, \sqrt[3]{2} \omega, \sqrt[3]{2} \omega^2) = \Q(\sqrt[3]{2}, \omega) = \sqrt(\sqrt[3]{2}, \sqrt{3} i) \text{ ya que } \omega = -\frac{1}{2} + \frac{\sqrt{3}i}{2}
        $$
        Además, como $\car(\Q) = 0$ es separable y $E = \Q(f)$ es normal, por tanto $E$ es de Galois.
        \item Hallamos el orden de la extensión.\\
        Podemos escribir el diagrama:\\
        \begin{center}
            \begin{tikzpicture}[node distance=1cm and 2cm, auto]
                \node (A) [above of =B] {$E = M(\omega)$};
                \node (B) {$M = \Q(\sqrt[3]{2})$};
                \node (C) [below of =B] {$\Q$};
                \draw[-] (A) to node {2 ($x^2+x+1$)}(B);
                \draw[-] (B) to node {3 ($x^3-2$)}(C);
            \end{tikzpicture}\\
        \end{center}
        Por tanto, por transitividad de grados $\vabs{E:\Q} = 6$ y una $\Q$-base de $E$ sería:
        $$
            \sdf{1, \sqrt[3]{2}, \sqrt[3]{4}, \omega, \sqrt[3]{2}\omega, \sqrt[3]{4}\omega}
        $$
        \item Determinar la clase de isomorfía de $G = \gal(E/\Q)$.\\
        Podemos verlo de dos formas:\\
        Sabemos que $\vabs{G} = \vabs{E:\Q} = 6$. Además por el teorema \ref{thm:3.11}, $G \leq S_3$ y por órdenes llegamos a que $G \isom S_3$.\\ % TODO: Explicar mejor (?)
        Por otra parte, sabemos que $M/\Q$ no es normal, por tanto por el TFTG (\ref{thm:fund-gal}), $\gal(E/M) \leq G$ no es normal, y entonces $G$ no es abeliano $\implies G \isom S_3$.

        \item Describir $G$ y sus órdenes.\\
        \begin{center}
            \begin{tikzpicture}[node distance=1cm and 2cm, auto]
                \node (A) [above of =B] {$E = L(\sqrt[3]{2})$};
                \node (B) {$L = \Q(\omega)$};
                \node (C) [below of =B] {$\Q$};
                \draw[-] (A) to node {$3$}(B);
                \draw[-] (B) to node {$2$}(C);
            \end{tikzpicture}\\
        \end{center}
        Como $L/\Q$ es normal (ya que es de grado $2$), todo $\tau \in G$ es la extensión de algún $\alpha \in \gal(L/\Q)$. Además, $\gal(L/\Q) = \sdf{1, \sigma}:\ \sigma(\omega) = \omega^2$.\\
        Cada elemento $\alpha \in \gal(L/\Q)$ se extiende a $E$ de tres formas distintas ($\tilde\alpha(\sqrt[3]{2}) \in \sdf{\sqrt[3]{2}, \sqrt[3]{2}\omega, \sqrt[3]{2}\omega^2}$) $\implies x^3-2 \in L[x]$ es irreducible y fijado por $\gal(L/\Q)$.
        Además, teniendo en cuenta que $\omega^2+\omega+1 = 0$ y $\omega^3 = 1$ podemos rellenar la siguiente tabla:
        \begin{center}
            \begin{tabular}{|c||c|c|c|}
            \hline
            Isomorfismo                 & $\omega$   & $\sqrt[3]{2}$         & ord\\ \hline \hline
            $\1_L   \rightarrow \tau_1 = 1_E$ & $\omega$   & $\sqrt[3]{2}$         & $1$\\ \hline
            $\1_L   \rightarrow \tau_2$       & $\omega$   & $\sqrt[3]{2}\omega$   & $3$\\ \hline
            $\1_L   \rightarrow \tau_3$       & $\omega$   & $\sqrt[3]{2}\omega^2$ & $3$\\ \hline
            $\sigma \rightarrow \tau_4$       & $\omega^2$ & $\sqrt[3]{2}$         & $2$\\ \hline
            $\sigma \rightarrow \tau_5$       & $\omega^2$ & $\sqrt[3]{2}\omega$   & $2$\\ \hline
            $\sigma \rightarrow \tau_6$       & $\omega^2$ & $\sqrt[3]{2}\omega^2$ & $2$\\ \hline
            \end{tabular}
        \end{center}
        Por los órdenes, y por ser un grupo de orden $6$ podemos determinar que es isomorfo a $S_3$ ya que no es cíclico y solo existen dos grupos de orden $6$. Otra forma sería intentar hallar una presentación del grupo a partir de relaciones entre los elementos, de esta forma hallaríamos la presentación:
        $$
            G = \gen{\tau_2 \tau_4 \mid o(\tau_2) = 3,\ o(\tau_4)=2,\ \tau_2^{\tau_4} = \tau_2^{-1}} = \gen{\tau_2} \ast \gen{\tau_4}
        $$
        \item Describir las subextensiones de $E/\Q$ indicando cuáles de ellas son normales.\\
        Por el TFTG:\\
        \begin{center}
            \begin{tikzpicture}[node distance=1cm and 2cm, auto]
                \node (B) {$K$};
                \node (A) [above of =B] {$\sdf{\Q \subseteq K \subseteq E}$};
                \node (E) [right= of B] {$\gal(E/K)$};
                \node (C) [below of =B] {$E^H$};
                \node (D) [above of =E] {$\sdf{H \leq G}$};
                \node (F) [below of =E] {$H$};
                \draw[<->] (A) to node {$1:1$}(D);
                \draw[|->] (B) to node {}(E);
                \draw[<-|] (F) to node {}(C);
            \end{tikzpicture}\\
        \end{center}
        $$
            K/\Q \text{ es normal } \iff \gal(E/K) \normsub G
        $$
        \begin{center}
            \begin{tikzpicture}[node distance=1cm and 2cm, auto]
                \node (A) {$E$};
                \node (B) [right= of A] {$\1$};
                \node (C) [below of =A] {$E^H$};
                \node (D) [below of =B] {$H$};
                \node (E) [below of =C] {$\Q$};
                \node (F) [below of =D] {$G$};
                \draw[-] (A) to node {}(C);
                \draw[-] (C) to node {$n$}(E);
                \draw[-] (B) to node {}(D);
                \draw[-] (D) to node {$n$}(F);
                \draw[<->] (A) to node {$1:1$}(B);
                \draw[<->] (E) to node {$1:1$}(F);
            \end{tikzpicture}\\
        \end{center}
        Como $\vabs{G} = 6$, si $1 < H < G$ ($H$ subgrupo propio), entonces por el teorema de Lagrange para grupos, si:
        \begin{align*}
            \vabs{H} = 2 \implies H \in Syl_2(G)\\
            \vabs{H} = 3 \implies H \in Syl_3(G)\\
        \end{align*}
        Entonces si $\vabs{H} = 3 \implies \vabs{G:H} = 2 \implies H \normsub G$. Con esto hemos visto que las subextensiones de $E/\Q$ son normales si son triviales o se corresponden con un grupo en $Syl_3(G)$.
        \begin{center}
            \begin{tikzpicture}[node distance=1cm and 2cm, auto]
                \node (A) {$P = \gen{\tau_2}$};
                \node (B) [right= of A] {$ $};
                \node (C) [right= of B] {$\gen{\tau_4} = Q$};
                \node (D) [above of =B] {$G$};
                \node (E) [below of =B] {$1$};
                \draw[=] (D) to node {}(A);
                \draw[-] (D) to node {$N_G(Q)$}(C);
                \draw[-] (A) to node {$3$}(E);
                \draw[-] (C) to node {$2$}(E);
            \end{tikzpicture}\\
        \end{center}
        Sabemos que $P \in Syl_3(G)$ y $Q \in Syl_2(G)$. Como $Q \centernot \normsub G$, consideremos el normalizador $N_G(Q) < G$, entonces $\vabs{N_G(Q):Q}\vabs{G:Q} = 3 \implies N_G(Q) = Q$. Por tanto $\vabs{Syl_2(G)} = 3$:
        $$
            Syl_2(G) = \sdf{Q,\ Q^{\tau_2},\ Q^{\tau_2^{-1}}} \text{ y } Q \neq Q^{\tau_2} \neq Q^{\tau_2^{-1}} \neq Q
        $$
        Y entonces $1, G, P=\gen{\tau_2}, Q, Q^{\tau_2}, Q^{\tau_3}$ son todos los subgrupos de $G$ donde $Q = \gen{\tau_4}$. Además, $E/\Q$ tiene una suextensión normal que es de grado $2$ sobre $\Q$. Sea $L$ esta extensión afirmamos que $L/\Q$ es normal y $\vabs{L:\Q} = 2$, $L=\Q(\omega)$ y por el TFTG: $L = E^{P} = E^{\gen{\tau_2}} = E^{\tau_2}$.

        Para calcular las subextensiones de grado $3$ sobre $\Q$, vamos a calcular $M_1 = E^{Q}$, $M_2 = E^{Q^{\tau_2}}$, $M_3 = E^{Q^{\tau_2^{-1}}}$. Además por el lema \ref{lm:4.5}:
        $$
            M_2 = E^{Q^{\tau_2}} = t_2(M_1),\ M_3 = E^{Q^{\tau_2^{-1}}} = \tau_2^{-1}(M_1)
        $$
        Vamos a calcular $M_1$, recordamos que $\sdf{1, \alpha, \alpha^2, \omega, \alpha \omega, \alpha^2 \omega}$ es una $\Q$-base de $E$ donde $\alpha = \sqrt[3]{2}$:
        \begin{align*}
            M_1 &= E^Q = E^{\gen{\tau_4}} = E^{\tau_4}\\
            E^{\tau^4} &= \sdf{x \in E:\ \tau_4(x) = x}\\
            x &= a + b\alpha + c \alpha^2 + d\omega + e\alpha\omega + f\alpha^2\omega = \tau_4(x) =\\
              &= \cancel a + \cancel{b\alpha} + \cancel{c\alpha^2} + d\omega^2 + e\alpha\omega^2 + f\alpha^2\omega^2 = (\text{ ya que } \omega^2+\omega+1 = 0)\\
              &= d(-\omega -1) + e\alpha^2(-\omega - 1) + f\alpha^2(-\omega -1) =\\
              &= d + 2d\omega + e\alpha + 2e\alpha \omega + f\alpha^2 + 2f\alpha^2\omega = 0 \implies\\
              &\implies d = e = f = 0
        \end{align*}
        Con este desarrollo hallamos:
        $$
            M_1 = E^{\tau_4} = \sdf{a + b\alpha + c\alpha^2 \mid a,b,c \in \Q} = \Q(\alpha) = \Q(\sqrt[3]{2})
        $$
        Para hallar $M_2$ y $M_3$ podemos repetir el mismo procedimiento o podemos:
        $$
            M_2 = E^{Q^{\tau_2}} = \tau_2(E^Q) = \tau_2(M_1) = Q(\sqrt[3]{2} \omega)
        $$
        y recordemos que $\tau_2^{-1} = \tau_3$, por tanto:
        $$
            M_3 = E^{Q^{\tau_3}} = \tau_3(E^Q) = \tau_3(M_1) = Q(\sqrt[3]{2} \omega^2)
        $$
    \end{enumerate}
\end{eg}

\begin{pro}[Consecuencia del teorema \ref{thm:3.11}]
    Sea $E/K$ una extensión, $E = K(f)$, $f \in K[x]$ irreducible y separable, y $\delta(f) = n$. Entonces $G = \gal(E/K) \leq S_n$ es transitivo, es decir:
    $$
        \forall i,j \in \sdf{1, \ldots, n} \exists \delta \in G \mid \delta(i) = j
    $$
\end{pro}
\begin{proof}
    Sea $\Omega = \sdf{\alpha_1, \ldots \alpha_n}$ el conjunto de raíces distintas de $f$ en $E$. Sea el morfismo $\rho: G \to S_\Omega;\ \delta \mapsto \rho(\delta)(\alpha) = \delta(\alpha)$, entonces como $f$ es separable $\ker(\rho) = \gal(E/E) = 1$ (tomando $L=E$ en \ref{thm:3.11}).\\

    Por tanto, $G \leq S_\Omega \isom S_n$, y como $f$ es irreducible, dados $\alpha_i, \alpha_j \in \Omega$ podemos afirmar que $\exists \delta \in G \mid \delta(\alpha_i) = \alpha_j$.
\end{proof}

\begin{obs}
    Los subgrupos transitivos de $S_n$ son $D_{2n}$
\end{obs}

\begin{pro}[Consecuencia del Teorema del Elemento Primitivo (\ref{thm:elemprim})]
    Sea $E/K$ una extensión de Galois, existe $f \in K[x]$ irreducible tal que $E = K(f)$.
\end{pro}

\begin{proof}
    Como $E/K$ es normal por ser de Galois, entonces $E = K(f)$, para algún $f \in K[x]$. Si $f$ es irreducible, lo hemos demostrado. Si $f$ no es irreducible, entonces $E = K_1(\alpha_1, \ldots, \alpha_n)$ y sabemos que por el teorema \ref{thm:elemprim}, $\exists \alpha \in E: E = K(\alpha)$. Entonces basta tomar $p = Irr(K, \alpha)$, con lo que tendríamos $E = K(p)$ con $p$ irreducible.
\end{proof}

\begin{pro}[Consecuencia del lema \ref{lm:4.5} y del teorema \ref{thm:4.3}]
    Sea $E/K$ una extensión de Galois, $G = \gal(E/K)$ su grupo de Galois, $M$ un cuerpo intermedio y $H = \gal(E/M) \leq G$, entonces:
    $$
        \forall \tau \in G,\ \tau(M) = E^{H^\tau}
    $$
\end{pro}
\begin{proof}
    Por el lema \ref{lm:4.5}, $H^\tau = \gal(E/\tau{M})$. Además, $E/K$ de Galois $\implies E/\tau(M)$ de Galois $\implies$ (por el teorema \ref{thm:4.3}) $E^{H^{\tau}} = \tau(M)$
\end{proof}

\begin{obs}
    El teorema \ref{thm:4.3} nos dice que sean $E/K$ finito y $G = \gal(E/K)$, si $E/K$ de Galois, entonces $E^G=K$. Se puede ver que el recíproco también es cierto. En resumen:\\

    \begin{center}
        Sean $E/K$ finita y $G=\gal(E/K)$:
        $$
            E^G = K \iff E/K \text{ es de Galois}
        $$
    \end{center}
\end{obs}

\begin{ex}[H4.11]
    Sea $E/K$ una extensión de Galois con $G = \gal(E/K)$ cíclico de orden $n$. Demuestra que:
    \begin{itemize}
        \item[(a)] Para cada divisor $d$ de $n$ existe exactamente un cuerpo intermedio $L$ con $\vabs{E:L} = d$.
        \item[(b)] Si $L_1$ y $L_2$ son dos cuerpos intermedios, entonces $L_1 \subseteq L_2$ si, y solo si, $E:L_2$ divide a $E:L_1$.
        \item[(c)] Demuestra que la hipótesis de $E/K$ de Galois puede relajarse a que $E/K$ sea finita.
    \end{itemize}

    \textbf{Solución}
    \begin{itemize}
        \item[(a)] Basta representar las correspondencias entre subcuerpos intermedios y subgrupos de Galois para comprobarlo.

        \begin{center}
            \begin{tikzpicture}[node distance=1cm and 2cm, auto]
                \node (A) {$E$};
                \node (B) [right= of A] {$\1$};
                \node (C) [below of =A] {$L = E^H$};
                \node (D) [below of =B] {$H$};
                \node (E) [below of =C] {$K$};
                \node (F) [below of =D] {$G$};
                \draw[-] (A) to node {$d$}(C);
                \draw[-] (C) to node {}(E);
                \draw[-] (B) to node {}(D);
                \draw[-] (D) to node {}(F);
                \draw[<->] (C) to node {$1:1$}(D);
            \end{tikzpicture}\\
        \end{center}
        \item[(b)]
        \begin{itemize}
            \item[$(\implies)$] Es trivial.
            \item[$(\Longleftarrow)$]
                \begin{align*}
                    d = &\vabs{E:L_2} = \vabs{H_2},\ H_2 = \gal(E/L_2)\\
                        &\vabs{E:L_1} = \vabs{H_1},\ H_1 = \gal(E/L_1) \implies d \divides \vabs{alpha_1} \divides n
                \end{align*}
                $H_1$ tiene un único subgrupo de orden $d$. Como $H_2$ es el único subgrupo de orden $d$ de $G \implies $ tienen que coincidir $\implies H_2 \subseteq H_1 \implies L_1 = E^{H_1} \subseteq L_2 = E^{H_2}$
        \end{itemize}
        \item[(c)]
            Basta aplicar que recíproco del teorema \ref{thm:4.3} (se resuelve en el ejercicio H4.4)
    \end{itemize}
\end{ex}

\begin{ex}[H4.8]
    Sea $E = \Q(\xi)$ donde $\xi = e^{\frac{2 \pi i}{7}}$. Muestra que $E$ es una extensión de Galois de $\Q$. Encuentra todos los subcuerpos intermedios de la extensión $E/\Q$, los subgrupos de $\gal(E/\Q)$ que les corresponden y determina qué subcuerpos intermedios se corresponden con extensiones normales de $\Q$.\\\\

    \textbf{Solución}\\
    $\xi$ es una raíz $p$-ésima de la unidad, con $p=7$. Es decir, $E = \Q(\xi) = \Q(x^7-1) = \Q(x^6 + x^5 + x^4 + x^3 + x^2 + x + 1)$. Por tanto, $\vabs{\Q(\xi):\Q} = 6$ y $G = \gal(\Q(\xi)/\Q) \isom \F_7^\times$ que es cíclico.\\
    $$G = \sdf{\varphi_j \mid 1 \leq j \leq 6} \mapsto \bar{j},\ \varphi(\xi) = \xi^j$$

    Buscamos $\gen{\varphi} = G$, que es lo mismo que buscar un generador de $\F_7^\times$ (de orden $6$). Fácilmente podemos ver que $3$ es un generador de $\F_7^\times$. Tomamos $\varphi = \varphi_3$, y como $G = C_6 \implies \varphi_2 = \varphi^2,\ \varphi_6 = \varphi^3$.\\
    Los subgrupos son:
    $$
        1,\ \gen{\varphi^2}, \gen{\varphi^3}, G
    $$
    Que se corresponden con las subextensiones (por el TFTG):
    $$
        E,\ L=E^{\varphi^2}, M=E^{\varphi^3}, \Q
    $$
    Vamos a calcular $L$:
    \begin{align*}
        L=E^{\varphi^2} &= \sdf{x = a_0 + a_1 \xi + \cdots a_5 \xi^5,\ a_i \in \Q \mid \varphi_2(x) = x}\\
        \varphi_2(x)    &= a_0 + a_1\xi^2 + a_2\xi^4 + a_3\xi^6 + a_4\xi + a_5\xi^3 =\\
                        &= a_0 + a_1\xi^2 + a_2\xi^4 + a_3(-1 -\xi \cdots -\xi^5) + a_5\xi^3 = x \iff \\
                        &\iff \xi(a_1-a_4-a_4) + \xi^2(a_2-a_1+a_3) + \xi^3(a_3-a_5+a_3) + \xi^4(a_4-a_2+a_3) + \xi^5(a_5-a_3) = 0
    \end{align*}
    %%TODO: REHACER FOTOS SANTORUM
    Con estas restricciones hallamos $L =\Q(\sqrt{-1})$.
\end{ex}

% % !TeX root = ../galois.tex

\chapter{Aplicaciones}

\section{Introducción. Radicales y resolubilidad.}

En este capítulo vamos a ver el \textit{Gran Torema de Galois}, que nos dice que un polinomio es \textit{resoluble por radicales} si y sólo si su grupo de Galois es \textit{resoluble}. Como es evidente, vamos a tener que introducir una serie de conceptos nuevos para desarrollar esta parte de la teoría.

\begin{dfn}[Radical]
    $R/K$ es \textbf{radical} si existe una cadena de subcuerpos $K = K_0 \subseteq K_1 \subseteq \cdots \subseteq K_{n-1} \subseteq K_n = R$, de modo que $K_j = K_{j-1}(a_j),\ (a_j)^{n_j} \in K_{j-1}$.

    \begin{center}
        \begin{tikzpicture}[node distance=1cm and 2cm, auto]
            \node (R)  {$R = K_n$};
            \node (Kj) [below of =R] {$K_j = K_{j-1}(a_j)$};
            \node (K2) [below of =Kj] {$K_2 = K_1(a_2)$};
            \node (K1) [below of =K2] {$a_2^{n_2} \in K_1 = K_0(a_1),\ x^{n_2}-a_2^{n_2}$};
            \node (K0) [below of =K1] {$a_1^{n_1} \in K = K_0 $};
            \draw[-] (R) to node {}(Kj);
            \draw[-] (Kj) to node {}(K2);
            \draw[-] (K2) to node {}(K1);
            \draw[-] (K1) to node {}(K0);
        \end{tikzpicture}\\
    \end{center}
\end{dfn}

\begin{eg}[Ejemplos de radicales]
    \begin{enumerate}[a)]
        \item Ejemplo 1:\\$$\Q(e^{\frac{2\pi i}{n}})/\Q,\ \ \ \ (e^{\frac{2\pi i}{n}})^n = 1 \in \Q$$
        \item Ejemplo 2:\\

        \begin{center}
            \begin{tikzpicture}[node distance=1cm and 2cm, auto]
                \node (A) {$ K_2 = \Q(\sqrt[3]{2}, i) = K_1(i) $};
                \node (B) [below of =A]{$ K_1 = \Q(\sqrt[3]{2}) = K_0(i)$};
                \node (C) [below of =B] {$ K_0 = \Q $};
                \draw[-] (A) to node {$i^2 = -1 \in K_1$}(B);
                \draw[-] (B) to node {$(\sqrt[3]{2})^3 = 2 \in K_0$}(C);
            \end{tikzpicture}\\
        \end{center}
        \item Ejemplo 3:\\

        \begin{center}
            \begin{tikzpicture}[node distance=1cm and 2cm, auto]
                \node (A) {$ \Q(\sqrt[3]{1+\sqrt 7})$};
                \node (B) [below of =A]{$ \Q(\sqrt 7)$};
                \node (C) [below of =B] {$ \Q $};
                \draw[-] (A) to node {$(\sqrt[3]{1+\sqrt 7})^3 = 1 + \sqrt{7} \in \Q(\sqrt 7)$}(B);
                \draw[-] (B) to node {$(\sqrt 7)^2 = 7 \in \Q $}(C);
            \end{tikzpicture}\\
        \end{center}
    \end{enumerate}
\end{eg}

\begin{dfn}[Resolubilidad por radicales]
    $f \in K[x]$ es \textbf{resoluble por radicales} si existe alguna extensión radical $R/K$ tal que $K(f) \subseteq R$.
\end{dfn}

\begin{eg}[Ejemplos de resolubilidad por radicales]
    \begin{enumerate}[a)]
        \item $x^p-1$. $\Q(x^p-1) = \Q(e^{\frac{2\pi i}{p}})$.
        \item $x^4-2$. $\Q(\sqrt[4]{2}, i)$.
        \item $x^4+4x^2+2$. $\Q(\sqrt{2+\sqrt{2}})$.
    \end{enumerate}
\end{eg}

\section{Grupos resolubles}
Como se anticipó en la sección anterior, el \textit{Gran Teorema de Galois} utiliza el concepto de grupo resoluble, para ello vamos a definir los conceptos necesarios.

\begin{dfn}[Grupo simple]
    Sea $G$ un grupo, es \textbf{simple} si no tiene subgrupos normales.
\end{dfn}

\begin{obs}$ $
    \begin{itemize}
        \item $G$ abeliano es simple si $G \isom C_p$.
        \item Sea $G_m$ el subgrupo máximo normal de $G$. Entonces $G_m \normsub G$ y no hay normales entre $G_m$ y $G \implies \quo{G}{G_m}$ es simple.
    \end{itemize}
\end{obs}

\begin{obs}[Clasificación de grupos simples]
    Si $G$ es finito y simple, entonces:
    \begin{itemize}
        \item $G \isom C_p$, si son abelianos.
        \item $G \isom A_n: n\geq5$.
        \item $G$ es de tipo \textit{Lie}.
        \item 26 esporádicos.
    \end{itemize}
\end{obs}

Los grupos finitos simples tienen relevancia en teoría de grupos ya que informalmente son los \textit{bloques} que forman el resto de grupos finitos. En cierta forma un grupo simple es a un grupo finito lo que un número primo a un número natural. De la misma forma que se puede descomponer en factores primos a un número natural, podemos escribir una descomposición de un grupo en subgrupos simples.

\begin{dfn}[Serie de descomposición]
    Sea $G$ un grupo finito, y $\sdf{H_j}_{j=1}^{j=n}$ una sucesión de subgrupos tales que:
    $$
        \1 = H_0 \normsub H_1 \normsub \cdots \normsub H_n = G
    $$
    se llama a dicha cadena \textbf{serie de descomposición} del grupo si $\forall j,\ \quo{H_{j+1}}{H_j}$ es un grupo simple y $n$ es la longitud de la serie.
\end{dfn}

\begin{obs}
    Aunque queda fuera del temario, todas las series de descomposición de un mismo grupo $G$ son equivalentes salvo por permutaciones e isomorfismos.
\end{obs}

\begin{dfn}[Grupo resoluble]
    Un grupo $G$ es \textbf{resoluble} si existe una serie de descomposición de la forma:
    $$
        \sdf{e} = G_0 \normsubeq G_1 \normsubeq G_2 \normsubeq \cdots \normsubeq G_{n-1} \normsubeq G_n = G
    $$
    tal que $\quo{G_j}{G_{j-1}}$ es abeliano.
\end{dfn}

\begin{pro}[Resolubilidad de grupos abelianos]
    Todo grupo $G$ abeliano es resoluble, es decir:
    $$
        G \text{ abeliano} \implies G \text{ resoluble.}
    $$
\end{pro}
\begin{proof}
    Como $\sdf{1} \subseteq G$, y $\sdf{1} \normsub G$ dado que $x \cdot 1 \cdot x^{-1} \in \sdf{1}$ y además $\quo{G}{\sdf{1}} \isom G$ es abeliano, hemos encontrado una serie:
    $$
        \sdf{1 = e} \normsub G
    $$
    tal que $\quo{G_j}{G_{j-1}}$ es abeliano.
\end{proof}
\begin{eg}$ $
    \begin{enumerate}[i)]
        \item $A_4$ es resoluble. Sea $V = \gen{(12)(34), (13)(24)}$:\\

        \begin{center}
            \begin{tikzpicture}[node distance=1cm and 2cm, auto]
                \node (A) {$ A_4 $};
                \node (B) [below of =A]{$ V $};
                \node (C) [below of =B] {$ \sdf{1} $};
                \draw[-, double] (A) to node {$3$}(B);
                \draw[-] (B) to node {$4$}(C);
            \end{tikzpicture}\\
        \end{center}
        \item $D_{2n}$ es resoluble.
        \item $S_3$ y $S_4$ son resolubles.
    \end{enumerate}
\end{eg}

\begin{obs}[Recuerdo de teoría de grupos]
    Sea $G$ un grupo, $N\normsubeq G$ y $H \leq G$ subgrupos y sea $M = H \cap N$, entonces:
    $$
        M \normsub H \text{ y } \quo{H}{M} \isom \quo{HN}{N}
    $$
    \begin{center}
        \begin{tikzpicture}[node distance=1cm and 2cm, auto]
            \node (A) {$N$};
            \node (B) [right= of A] {$ $};
            \node (C) [right= of B] {$H$};
            \node (D) [above of =B] {$G$};
            \node (E) [below of =B] {$N\cap H = M$};
            \node (F) [below of =E] {$\1$};
            \draw[-, double] (D) to node {}(A);
            \draw[-] (D) to node {}(C);
            \draw[-] (A) to node {}(E);
            \draw[-] (C) to node {}(E);
            \draw[-] (E) to node {}(F);
        \end{tikzpicture}\\
    \end{center}
\end{obs}

\section{Resolubilidad y extensiones de Galois}

En esta sección vamos a ver resultados que relacionen los conceptos de un radical $R/K$ con extensiones de Galois, ya que en general un radical no es una extensión de Galois, como por ejemplo $\Q(\sqrt[4]{2})/\Q$ y
$\Q(\sqrt{1+\sqrt 7})/\Q$.

\begin{thm}[Kummer]\label{thm:5.1}
    Sean $E, K$ cuerpos de la forma $\Q \subseteq K \subseteq E \subseteq \C$, si $K$ contiene $\xi \in \C$ con $o(\xi) = n$, entonces son equivalentes:
    \begin{enumerate}[a)]
        \item $E = K(a)$ con $a \in E$ tal que $a^n \in K$.
        \item $E/K$ es una extensión de Galois con $\gal(E/K)$ cíclico de orden un divisor de $n$.
    \end{enumerate}
\end{thm}

\begin{proof}$ $
    \begin{enumerate}
        \item[a) $\implies$ b)] Tenemos que $E = K(a)$ con $a^n \in K$. Podemos suponer que $a \notin \K$. Sea $f(x) = x^n - a^n \in K[x]$, las raíces de $f$ son $\sdf{a\xi^j \mid 0 \leq j \leq n}$.
        Como $E = K(a) = K(a\xi^j\mid 0 \leq j \leq n) = K(x^n - a^n)$, entonces $E/K$ es de Galois.\\
        Sea $\sigma \in \gal(E/K)$, $\sigma(a) = a\xi^{\kappa_\sigma}$ donde $0 \leq \kappa_\sigma \leq n$. Definimos:\\
        \begin{center}
            \begin{tikzpicture}[node distance=1cm and 2cm, auto]
                \node (B)                {$\sigma$};
                \node (A)  [above of =B] {$\rho: \gal(E/K)$};
                \node (E)  [right= of B] {$\xi^{\kappa_\sigma}$};
                \node (D)  [above of =E] {$\gen{\xi}$};
                \draw[->]  (A) to node {}(D);
                \draw[|->] (B) to node {}(E);
            \end{tikzpicture}\\
        \end{center}
        Y entonces:
        \begin{itemize}
            \item $\rho(\1_E) = \rho^0 = 1$
            \item $\rho(\sigma\tau) = \xi^{\kappa_\sigma+\kappa_\tau} \xi^{\kappa_\sigma} * \xi^{\kappa_\tau} = \rho(\sigma)\rho(\tau)$.
            Donde $\sigma\tau(a) = \tau(a \xi^{\kappa_\sigma}) = a \xi^{\kappa_\tau} \xi^{\kappa_\sigma} = \xi^{\kappa_\sigma+\kappa_\tau}$
            \item $\rho$ es inyectiva porque $\rho(\sigma) = 1 \implies \sigma(a) = a \implies \sigma = 1_E$ y por tanto $\gal(E/K) \leq \gen{\xi}$.
        \end{itemize}
        \item[b) $\implies$ a)] Como $E/K$ es de Galois, $\gal(E/K) = \gen{\tau}$ y sea $n$ el orden del grupo de Galois, entonces: $o(\tau) = d\divides n$ y $\vabs{E:K} = d$.

        Sea $\omega \in\gen{\xi}$ tal que $o(\omega) = d$ (por ejemplo $\omega = \xi^{\frac{n}{d}})$, basta encontrar un $a \in E$ tal que $\tau(a) = a\omega$ y $a \notin K$. Vamos a ver por qué:\\
        Supongamos que existe tal $a$, entonces:
        $$
            \tau(a^d) = (\tau(a))^d = a^d \omega^d = a^d \implies a^d \in E^\gen{\tau} = K.
        $$
        En particular, $a^n = (a^d)^{\frac{n}{d}} \in K$. Sea $p = Irr(K, a) \in K[x]$, $p \divides x^d - a^d$.
        Por otro lado, $a, a\omega = \tau(a), \ldots, a\omega^{d-1} = \tau^{d-1}(a)$ son raíces de $p$ distintas $\implies \delta(p) \geq d \implies p(x) = x^d - a^d$, por tanto:
        $$
            \vabs{K(a):K} = \delta(p) = d = \vabs{E:K} \implies E = K(a)
        $$

        Finalmente, solo nos queda buscar el $a$ que mencionamos arriba. Por el \hyperref[thm:4.2]{teorema de Dedekind}, $\sdf{1, \tau, \ldots, \tau^{d-1}}$ es linealmente independiente y entonces:
        $$
            \1 + \frac{1}{\omega} \tau + \cdots + \frac{1}{\omega^{d-1}}\tau^{d-1} \neq 0
        $$
        Sea $b \in E$ tal que $0 \neq b + \frac{1}{\omega} \tau(b) + \cdots + \frac{1}{\omega^{d-1}}\tau^{d-1}(b) = a$ y hemos encontrado el $a$ que buscábamos. Puede comprobarse que $\tau(a) = \omega a = a \omega$.
    \end{enumerate}
\end{proof}

El siguiente grafo resume el resultado del teorema:\\
\begin{center}
    \begin{tikzpicture}[edge quotes mid, node distance=1cm and 2cm, auto]
        \node (A) [above of =C]{$E(\xi) = \gen{K(\xi), E}$};
        \node (B) [] {$K(\xi)$};
        \node (C) [right= of B] {$ $};
        \node (D) [right= of C] {$E = K(a)$};
        \node (E) [below of =C] {$K(\xi)\cap E$};
        \node (F) [below of =E] {$K$};
        \draw[-] (A) to node [sloped]{abeliana}(B);
        \draw[-] (A) to node {}(D);
        \draw[-] (B) to node {}(E);
        \draw[-] (D) to node {}(E);
        \draw[-] (B) to node [sloped, below]{abeliana}(F);
    \end{tikzpicture}\\
\end{center}

\begin{ex}[H5.2]
    Sea $E/K$ una extensión de Galois, $K \subseteq L, M \subseteq E$ subcuerpos intermedios. $\gen{L, M} = \bigcap_{H \subseteq E} H$ tal que   $L,M \subseteq H$. Prueba que:
    $$
        \gal(E/L) \cap \gal(E/M) = \gal(E/\gen{L, M})
    $$
    \textbf{Solución}\\\\

    Vamos a resolverlo por doble inclusión. Por definición $\gal(E/\gen{L, M}) \subseteq \gal(E/L) \cap \gal(E/M)$.\\
    Además $L,M\subseteq E^{\gal(E/L) \cap \gal(E/M)} \subseteq E$ y $\gen{L, M} \subseteq E^{\gal(E/L) \cap \gal(E/M)}$. Por tanto:
    $$
        \gal(E/E^{\gal(E/L) \cap \gal(E/M)}) \subseteq \gal(E/\gen{L, M})
    $$
    y $\gal(E/L) \cap \gal(E/M) = \gal(E/E^{\gal(E/L) \cap \gal(E/M)})$ por el \hyperref[thm:fund-gal]{TFTG}.
\end{ex}

\begin{ex}[H5.4]
    Por completar. %TODO
\end{ex}

\section{Gran Teorema de Galois}

\textit{NOTA: Solo están redactados los resultados teóricos de esta sección. Todos los ejercicios de esta sección están por redactar (se corresponden con H5.8, H5.10, H5.11, H5.12, H5.5, H5.6). También falta algún ejemplo, y las pruebas de los resultados.}

\begin{obs}[Recuerdo de teoría de grupos]
    Sea $G$ un grupo y $N \normsubeq G$, $H \leq G$ subgrupos, entonces:
    \begin{enumerate}[a)]
        \item $NH \leq G$, $N \cap H \normsubeq H$ y además $$\quo{NH}{N} \isom \quo{H}{N\cap H}$$.
        \item Si además $H \normsub G$, entonces $NH \normsub G$, $H  \normsubeq NH$ y además:
        $$
            \quo{G}{NH} \isom \quo{\quo{G}{N}}{\quo{NH}{N}}
        $$
    \end{enumerate}
\end{obs}

\begin{dfn}[Extensión abeliana]
    Una extensión $E/K$ es \textbf{abeliana} si $E/K$ es de Galois y $\gal(E/K)$ es un grupo abeliano.
\end{dfn}

\begin{lm}[Extensión de Galois y resoluble]\label{lm:5.2}
    Sea $\Q \subseteq K$, $L \subseteq \C$ y $L/K$ una extensión. Supongamos que existe una sucesión $\sdf{K_j}_{j=0}^{j=m}$ tal que $K = K_0$, $L = K_m$ monótona creciente por inclusión y $K_j/K_{j-1}$ es abeliana. Entonces si $E/K$ es de Galois, con $E \subseteq L$, entonces $\gal(E/K)$ es resoluble.
\end{lm}
\begin{proof}
    A completar.
\end{proof}
\begin{thm}[Gran Teorema de Galois]\label{thm:5.3}\label{thm:GTG}
    Sea $\Q \subseteq K$ un cuerpo, $f \in K[x]$ un polinomio y $E = K(f) \subseteq \C$ su cuerpo de escisión.
    $$
        f \text{ es resoluble por radicales } \iff \gal(f) = \gal(E/K) \text{ es resoluble }
    $$
\end{thm}
\begin{proof}
    A completar.
\end{proof}

\begin{lm}[Generación del grupo simétrico]\label{lm:5.4}
    Para $n \leq 2$, $S_n = \gen{(12)(12\cdots n)}$
\end{lm}

\begin{thm}[de Abel-Ruffini]
    Las ecuaciones polinómicas \textbf{generales} $f(x) = a_n x^n + \cdots + a_1 x + a_0 = 0$ no son resolubles por radicales cuando $\delta(f)\geq 5$. Es decir, no es posible encontrar las soluciones de $f$ de forma general aplicando finitamente la suma, resta, multiplicación y radicación.
\end{thm}

\begin{proof}
    A completar.
\end{proof}

\begin{obs}
    Equivalentemente, el teorema enuncia que no existe una fórmula general para hallar las raíces de $f(x)$ con $\delta(f) \geq 5$ aplicada a sus coeficientes. Para $\delta(f) \leq 4$ si es posible, por ejemplo con $\delta(f) = 2$ tenemos:
    $$
        x = \frac{-b \pm \sqrt{b^2-4ac}}{2a}
    $$
    Existen también fórmulas para polinomios de tercer y cuarto grado.
\end{obs}

\part{Apéndices}

%\renewcommand{\chaptername}{Apéndice}
%\renewcommand{\thechapter}{\Alph{chapter}}
%\setcounter{chapter}{0}

%\input{input/a-ejercicios}

%\renewcommand{\chaptername}{Apéndice}
%\renewcommand{\thechapter}{\Alph{chapter}}
%\setcounter{chapter}{0}

%\input{input/a-ejercicios}

\chapter{\'{I}ndices}

%%TODO por qué no podemos utilizar only= para hacer whitelisting en lugar de blacklisting? que alguien lo arregle parfavar
\renewcommand{\listtheoremname}{Lista de definiciones}
\listoftheorems[ignore={thm,eg,pro,cor,obs,lm,ex,th_ex}]

\renewcommand{\listtheoremname}{Lista de teoremas}
\listoftheorems[onlynamed,ignore={dfn,eg,cor,obs,ex,th_ex}]

\renewcommand{\listtheoremname}{Lista de ejemplos}
\listoftheorems[onlynamed,ignore={dfn,thm,pro,cor,obs,lm,ex,th_ex}]

\renewcommand{\listtheoremname}{Lista de ejercicios}
\listoftheorems[ignore={dfn,thm,pro,cor,obs,lm,eg}]

\printglossaries

\bibliographystyle{alpha}

\end{document}
